\documentclass[12pt]{article}


% Basic packages only
\usepackage[a4paper, margin=20mm]{geometry}
\usepackage{amsmath}
\usepackage{enumitem}
\usepackage{graphicx}
\usepackage{microtype}
\usepackage{longtable}
\usepackage{array}
\usepackage{fontspec}
\usepackage{tfrupee} 
\setmainfont{Latin Modern Roman}

% Better font and spacing
\renewcommand{\familydefault}{\rmdefault}
\setlength{\parindent}{0pt}
\setlength{\parskip}{8pt}
\setlist{nosep, leftmargin=18pt}

% Section formatting
\usepackage{titlesec}
\titleformat{\section}{\large\bfseries}{\thesection}{1em}{}
\titleformat{\subsection}{\normalsize\bfseries}{\thesubsection}{1em}{}

\begin{document}

\begin{center}
{\Large \textbf{Stock Analysis}}\\[0.5em]
{\large Fundamental Analysis and Financial Modelling}
\end{center}

\section{How Share Prices Move/Fluctuate}

Share prices are determined by the basic economic principle of supply and demand:

\begin{center}
\textit{Demand $\uparrow$ Supply $\downarrow$ $\implies$ Price $\uparrow$}\\
\textit{Demand $\downarrow$ Supply $\uparrow$ $\implies$ Price $\downarrow$}\\[0.5em]
\textbf{Key Driver: Information}
\end{center}

\subsection{Information Sources}

Market information comes from various channels that influence investor decisions:

\begin{itemize}
\item Annual Reports
\item Quarterly Earnings
\item Court Cases
\item Press Releases
\item News Stories
\item Social Media
\item Market Buzz
\end{itemize}

\subsection{Primary Influencing Factors}

Share price movements are influenced by four major categories of factors:

\textbf{1. Company Fundamentals}

The fundamental strength of a company forms the foundation of its valuation:
\begin{itemize}
\item Earnings
\item Growth Prospects
\item Management \& Governance
\item Competitive Advantage
\item Expansion Plans
\item Debt Position
\item Promoter Shareholding
\item Pledging of Shares
\end{itemize}

\textbf{2. Technical Indicators}

Technical factors reflect market sentiment and trading dynamics:
\begin{itemize}
\item Liquidity
\item Positive/Negative News Flow
\item Chart Patterns
\item Market Valuation (Overpriced/Underpriced)
\end{itemize}

\textbf{3. Micro \& Macro Economic Factors}

Broader economic conditions significantly impact stock performance:
\begin{itemize}
\item Socio-Political Stability/Instability
\item Government Policies
\item National \& Global Economy/GDP Trends
\item Inflation Rates
\item Competitive Investment Landscape
\end{itemize}

\textbf{4. Human Sentiment}

Market psychology plays a crucial role in price movements:
\begin{itemize}
\item Information Asymmetry
\item Divergent Reactions to a Singular Piece of Information
\end{itemize}

\begin{center}
\textit{(Trend is your friend. Follow the trend, not news.)}
\end{center}

\subsection{Investing vs Trading}

The two primary approaches to market participation differ significantly:

\begin{center}
\begin{tabular}{|>{\raggedright}p{6.5cm}|>{\raggedright\arraybackslash}p{6.5cm}|}
\hline
\textbf{Investing} & \textbf{Trading} \\
\hline
Generally for the long term & Typically for the short or very short term \\
\hline
Utilizes Fundamental Analysis & Relies on Technical Analysis \\
\hline
Associated with lower risk & Involves higher risk \\
\hline
Aims for wealth creation & Aims to generate profits (make money) \\
\hline
Requires less active time participation & Demands more active time commitment \\
\hline
\end{tabular}
\end{center}

\section{Investment and Trading Strategies}

\subsection{Types of Investing}

\textbf{Growth Investing}

Growth investors focus on companies with exceptional growth potential:
\begin{itemize}
\item Growth investors are attracted to companies that are expected to grow faster (by revenues/cash flow and profits) than the rest.
\item \textit{Father: Philip Fisher}
\end{itemize}

\textbf{Value Investing}

Value investors seek undervalued opportunities in the market:
\begin{itemize}
\item It is about finding companies whose stock price doesn't necessarily reflect their fundamental worth.
\item \textit{Father: Benjamin Graham}
\end{itemize}

\subsection{Types of Trading}

Different trading strategies are categorized by their time horizons:

\begin{enumerate}
\item \textbf{Scalp Trading}
    \begin{itemize}
    \item Duration: Seconds, minutes; Buy-Sell
    \end{itemize}

\item \textbf{Intraday Trading}
    \begin{itemize}
    \item Duration: 1 day; Buy-Sell
    \end{itemize}

\item \textbf{Swing Trading}
    \begin{itemize}
    \item Duration: 1-7 days; based on short-term trends
    \end{itemize}

\item \textbf{Momentum Trading}
    \begin{itemize}
    \item Strategy: Buy securities that are rising and sell when they appear to have peaked
    \end{itemize}

\item \textbf{Positional Trading}
    \begin{itemize}
    \item Duration: Months; based on long-term potential of stock
    \end{itemize}

\item \textbf{BTST (Buy Today Sell Tomorrow) \& STBT (Sell Today Buy Tomorrow)}
\end{enumerate}

\subsection{Analysis Methods}

\textbf{Fundamental Analysis}

Fundamental analysis provides the foundation for investment decisions:
\begin{itemize}
\item Attempts to evaluate stocks by measuring their intrinsic value.
\item Involves studying everything from the overall economy and industry to the financial strength and management of individual companies.
\item The assumption is that in the longer horizon, the share price of a company with good fundamentals will rise, thus providing profits to investors.
\item Good for long-term investments.
\item Generally less risky.
\item Can be more time-consuming.
\end{itemize}

\textbf{Recommended Books:}
\begin{enumerate}
\item The Intelligent Investor
\item Security Analysis
\item Common Stocks and Uncommon Profits
\end{enumerate}

\textbf{Technical Analysis}

Technical analysis serves as the primary tool for trading decisions:
\begin{itemize}
\item Attempts to identify trading opportunities by looking at statistical trends, such as movements in a stock's price and volume.
\item The assumption is that all known fundamentals are already factored into the price and that history tends to repeat itself.
\item Primarily used for trading.
\item Generally involves higher risk.
\item Can be less time-consuming compared to fundamental analysis.
\end{itemize}

\textbf{Recommended Books:}
\begin{enumerate}
\item Technical Analysis of the Financial Markets by John J. Murphy
\item Numerous resources available, including YouTube videos for learning Technical Analysis
\end{enumerate}

\section{Fundamental Analysis}
Fundamental analysis involves evaluating a company's intrinsic value through three layers of examination:

\subsection{Economy Analysis}
When people make transcations (buy-sell product and services) it affects the economy of a country. We need to perform economy analysis to get an idea for right time and industry to invest in.

\begin{center}
\begin{figure}[!h]
\centering
\includegraphics[width=0.5\textwidth]{images/economy.png}
\caption{Economic movementa}
\label{img:anything}
\end{figure}
\end{center}

Understanding macroeconomic conditions that impact investment decisions:

\begin{itemize}
\item \textbf{Crude Oil Prices}
    \begin{itemize}
    \item Increase $\rightarrow$ Economic Health $\downarrow$, Inflation $\uparrow$
    \item Decrease $\rightarrow$ Good Economic Health $\leftrightarrow$ (for India)
    \item Impact: For every \$10 increase $\rightarrow$ GDP decreases by 0.4\%
    \item Affected sectors: Exploration \& Production, Refining \& Distribution, Lubricants, Chemicals, Diesel, Petrol, Natural gas, Jet fuel, LPG, Waxes, Bitumen for roads and roofing, Polishes, fuels for ships, factories like Plastic, medicine, rubber, etc.
    \end{itemize}

\item \textbf{GDP (Gross Domestic Product)}
    \begin{itemize}
    \item Represents total production value of a nation
    \item Higher values generally better, but lower GDP may present buying opportunities
    \end{itemize}

\item \textbf{Inflation}
    \begin{itemize}
    \item Measured by Consumer Price Index (CPI) \& Wholesale Price Index (WPI)
    \item Lower inflation preferred for economic stability
    \end{itemize}

\item \textbf{Monsoons}
    \begin{itemize}
    \item 4 months in a year accounts for 75\% of annual rainfall in India.
    \item Affects Agriculture sector which accounts for 15\% of GDP, 58\% employment and 70\% of dependency for Indian Economy.
    \item Affected sectors: Agriculture inputs (seeds, agrochemicals, fertilizers, etc.), Auto, FMCG
    \item In April or 1st week of June, monsoon announcement is made by IMD (Indian Meteorological Department) which is a key indicator for the market especially the sectors affected by it directly.
    \item Better monsoon $\rightarrow$ Better agriculture output $\rightarrow$ Better GDP $\rightarrow$ Better market performance (but floods are not good)
    \end{itemize}

\item \textbf{Elections}
    \begin{itemize}
    \item Political stability impacts market confidence
    \item Election Results that are too variable can lead to market volatility
    \item Stable governments generally lead to better market performance
    \item Policy changes can create sector-specific opportunities
    \end{itemize}

\item \textbf{Indian Rupee Valuation}
    \begin{itemize}
    \item Strong rupee benefits importers, weak rupee helps exporters
    \item Rising value of rupee leads to better economic health, better fiscal deficit, and lower inflation, good interest rate on FDI (Foreign Direct Investment) Loans, also makes import/export easier and reasonable.
    \item Key factors: Trade wars, interest rates, geopolitical tensions, import-export, economic downturn, unforeseen election results, shaky political reforms
    \item Sensitive sectors: With rising value of rupee, sectors like Pharma, IT, and Jewelry may face challenges due to reduced competitiveness in exports. While sectors like Oil and Gas, Cement, Tiles, FMCG, Plastic, Chemicals, Tyre, etc. which import raw materials, may benefit from a stronger rupee.
    \end{itemize}
\end{itemize}

Other economic indicators:
\begin{itemize}
\item Forex reserves \& trade balance
\item Agricultural output \& tax structures
\item Interest rates \& purchasing power parity
\item Government infrastructure spending
\item Tax structures \& fiscal policies
\item Foreign Direct Investment (FDI) and FII (Foreign Institutional Investment) trends.
\item Per Capita Income \& purchasing Power Parity (PPP)
\item Employment rates \& wage growth
\end{itemize}

\subsection{Industry Analysis}
We need to perform industry analysis to get an idea about what industry will be the best to invest in based on its future prospects (more growth potential, high profitability and lower risks). Evaluating sector-specific opportunities and risks:
\begin{itemize}
\item Growth potential vs competitive landscape
\item Government policies \& regulatory environment
\item Technological disruption \& substitution threats
\item Raw material availability \& input costs
\end{itemize}


\[
\text{Stock Market growth} ~ \alpha ~ \text{Industry potential}
\]
Industry potential depends on its growth, risks, opportunity and competition.

\subsection{Factors affecting Industry}


\begin{itemize}
\item \textbf{Industry Life Cycle Stages}:
    \begin{enumerate}
    \item Pioneering Stage (Emerging technologies)
    \item Rapid Growth Stage
    \item Maturity \& Stabilization
    \item Decline/Diversification
    \end{enumerate}
    
\item \textbf{Competitive Landscape}:
    \begin{itemize}
    \item Number of competitors and market share distribution
    \item Barriers to entry (capital requirements, regulations)
    \item Product differentiation level
    \end{itemize}

\item \textbf{Government Influence}:
    \begin{itemize}
    \item Policy support or restrictions
    \item Tax incentives/subsidies
    \item Environmental regulations
    \end{itemize}

\item \textbf{Supply Chain Dynamics}:
    \begin{itemize}
    \item Raw material availability
    \item Supplier bargaining power
    \[
    \text{Supplier Power} = \frac{\text{Number of Suppliers}}{\text{Switching Costs}}
    \]
    \end{itemize}

\item \textbf{Other Qualitative Factors}:
    \begin{itemize}
    \item Labour availability and skill levels required
    \item Technological advancements and innovation pace
    \item Consumer preferences and demand trends. Business cycles and seasonality.
    \item Threats from substitutes and new entrants
    \item Globalization impact and international competition
    \end{itemize}

\item \textbf{Quantitative Factors}:
    \begin{itemize}
    \item Market size and growth rate
    \item Profit margins and return on investment (ROI)
    \item Turnover and Profitability (Sales) trends
    \item Cost Structure (fixed vs variable costs, operating leverage, cost competitiveness within the industry, scalability of operations, and sensitivity to input price changes)
    \end{itemize}

\item \textbf{Other Factors}:
    \begin{itemize}
        \item Foreign trade and investments
        \item Industrial Production
        \item International Events (e.g., discovery of precious metal deposits, new technology breakthroughs, etc.)
    \end{itemize}
\end{itemize}

\section{Company Analysis}
Fundamental analysis of individual companies involves both qualitative and quantitative evaluation:

\subsection{Qualitative Analysis}
\subsubsection*{(i) Business Model Evaluation}
Bussiness model of a company should be simply explainable to even a children. Key questions to assess business viability:
\begin{enumerate}
\item What does the company do? (Core operations)
\item Manufacturing capabilities \& service delivery mechanisms
\item Plant locations \& capacity utilization
\item Raw material sourcing \& input costs
\item Client portfolio \& end-user demographics
\item Competitive landscape analysis
\item Geographic expansion plans
\item New product pipeline
\item Revenue mix breakdown (revenue model and sources)
\item Best selling Products/Services
\item Distribution channels \& sales strategies
\item Regulatory environment constraints
\item Banking partners \& audit quality
\item Workforce composition \& skill levels
\item Corporate structure (subsidiaries/associates/alliances/parent company)
\item Management type (Professional vs Family-run)
\end{enumerate}

\subsection*{$\rightarrow$ SWOT Analysis}

SWOT Analysis is a strategic tool used to evaluate the internal and external factors that affect a business. It stands for:

\begin{itemize}
  \item \textbf{Strengths:} Internal attributes that support successful outcomes.
    \begin{itemize}
        \item Strong management
        \item Quality products and services
        \item High profit margins
        \item Loyal customer base
        \item Sound financial position
        \item Cost advantage
    \end{itemize}

  \item \textbf{Weaknesses:} Internal limitations or challenges that hinder performance.
    \begin{itemize}
        \item Lack of funds
        \item Product liability risks
        \item Absence of innovation
        \item Unstable management
        \item Competitive intensity
    \end{itemize}

  \item \textbf{Opportunities:} External conditions that the business can capitalize on.
    \begin{itemize}
        \item Internal and external growth potential
        \item Horizontal, vertical, and technical expansion
        \item Emerging social trends
        \item Favorable macroeconomic conditions
    \end{itemize}

  \item \textbf{Threats:} External elements that could challenge business performance.
    \begin{itemize}
        \item Litigation risks
        \item Changing government policies
        \item Intense competition
        \item Negative market sentiment
        \item Substitutes and pricing pressure
    \end{itemize}
\end{itemize}

\vspace{5pt}
\noindent SWOT can be used to evaluate not only businesses, but also industries at large, by identifying the key drivers and potential risks in the sector.


\subsubsection*{(ii) Competitive Advantage (Economic Moat)}

According to Warren Buffett, a company with a sustainable competitive advantage is one that can maintain its profitability and market position over the long term. This is often referred to as having an "economic moat."
An economic moat is a unique advantage that protects a company from its competitors, allowing it to maintain higher profit margins and market share.
\\ Characteristics of companies with sustainable advantages:
\begin{itemize}
\item High customer switching costs
\item Strong network effects
\item Operational cost leadership (low operation costs)
\item Proprietary assets/IP
\item Pricing power (near-monopoly status)
\end{itemize}

\textbf{Identification Parameters:}
\begin{enumerate}
\item Substantial cash reserves
\item Consistent performance during economic downturns
\item Brand equity \& R\&D capability
\item Product line diversification
\item Large-scale operations
\item Regulatory protections (Goverment and regulatory licenses)
\item High customer retention rates
\end{enumerate}

\subsection*{$\rightarrow$ Porter's Five Forces (Company Perspective)}
Porter's Five Forces is a framework for analyzing the competitive intensity and attractiveness of an industry/company. It helps assess a company's strategic position by examining five key forces:

\begin{enumerate}
  \item \textbf{Rivalry Among Existing Competitors}
  \begin{itemize}
    \item Number of competitors in the market
    \item Market exit barriers
    \item Industry size and growth rate
    \item Degree of product differentiation or homogeneity
    \item Size and loyalty of customer base
    \item Threat of horizontal expansion
  \end{itemize}

  \item \textbf{Threat of New Entrants}
  \begin{itemize}
    \item Capital required to start a new business
    \item Legal and regulatory barriers to entry
    \item Access to supplier and distribution networks
    \item How existing firms react to new entrants
  \end{itemize}

  \item \textbf{Threat of Substitute Products or Services}
  \begin{itemize}
    \item Number and availability of substitutes
    \item Performance and quality of substitutes
    \item Cost incurred by customers when switching to substitutes
  \end{itemize}

  \item \textbf{Bargaining Power of Buyers}
  \begin{itemize}
    \item Number of buyers
    \item Size and value of individual purchases
    \item Switching costs faced by buyers
    \item Price sensitivity of buyers
  \end{itemize}

  \item \textbf{Bargaining Power of Suppliers}
  \begin{itemize}
    \item Number of suppliers
    \item Financial strength of suppliers
    \item Cost of switching to alternative suppliers
    \item Availability of substitute raw materials or inputs
  \end{itemize}
\end{enumerate}

\subsubsection*{(iii) Management Analysis}
Company's management are the people whom you are trusting with your money and expect that they will generate some returns for you, thus management analysis plays an integral role in your investment's success.
Understanding the quality and intent of company management is crucial for investment decisions. Key considerations include:

\begin{enumerate}
    \item \textbf{Promoter Background} \\
    Review track records for frauds, regulatory issues (SEBI, RBI), and ongoing legal cases.
  
    \item \textbf{Remuneration of Senior Management} \\
    Analyze salary trends, allowances, personal warrent issues, and perks relative to company profits (e.g., year-over-year increase).
  
    \item \textbf{Related Party Transactions} \\
    Scrutinize dealings with promoter-linked entities, joint ventures, interest free loans to relatives, two or more same businesses by same management, or personal loans on behalf of the company.
  
    \item \textbf{Management Forecast} \\
    Evaluate how past projections fared (to which extent the management acieved past projections) and whether recent strategic changes were successfully implemented. Future plans.
  
    \item \textbf{Dividend Policy} \\
    Understand how profits are distributed whether management rewards shareholders or retains earnings. Check the source of money for dividend payments.

    \item \textbf{Shareholding Patterns and Pledging} \\
    This analysis should cover the following key areas:
    \begin{itemize}
      \item Check FII/DII holdings, promoter stake changes, and pledging trends. High pledging implies risk. 
      \item Who are significant shareholders? Are they increasing or decreasing their stake?
      \item Have Mutual Funds and Investment Firms like Insurance companies and Banks invested in the company?
      \item Are they increasing or decreasing their stake? (A positive trend shows confidence in the company)
      \item Note: Professionally-run businesses (e.g., by banks) with transparent governance are generally preferred.
      \item High promoter pledging or opaque accounting may signal financial instability or value traps.
      \item But it can also be a sign of confidence if the company is using the funds for growth and expansion (check if company's cash flow is increasing).
    \end{itemize}
  
    \item \textbf{Account Manipulation} \\
    Watch for complex or inconsistent accounting practices that mask the company's real performance.
\end{enumerate}


\subsubsection*{(iv) Corporate Governance}

Corporate governance defines the rules and mechanisms for managing a company in the interest of shareholders. Important elements include:

\begin{enumerate}
  \item \textbf{Ownership and Voting Control} \\
    Analyze promoter holding, voting rights, and control structures.
  
  \item \textbf{Board of Directors Representation} \\
    Assess board independence and whether it meets current and future needs.
  
  \item \textbf{Remuneration Linked to Performance} \\
    Executive pay should correlate with company performance.
  
  \item \textbf{Investor Impact} \\
    Understand the role and influence of institutional investors.
  
  \item \textbf{Shareholder Rights} \\
    Ensure minority shareholders have adequate protection and influence.
  
  \item \textbf{Risk Management} \\
    Examine how long-term risks are addressed through policies and strategy.
\end{enumerate}

\textbf{Red Flags:} Poor governance undermines transparency, financial reliability, and shareholder trust.

\begin{itemize}
  \item Are company rules aligned with mission and vision?
  \item Is every stakeholder served fairly?
  \item Is the company legally compliant with policies?
\end{itemize}

\vspace{0.5em}
\textbf{Other Qualitative Factors:}
\begin{itemize}
  \item Customer/geographic exposure, market share, disruptive tech
  \item Regulatory issues, share transactions, management vision
  \item Reasons for Corporate Issues (e.g., Dividents Payout, Stock Split, Bonus Shares, Rights Issue)
  \item Corporate culture, employee satisfaction, and retention
  \item Corporate social responsibility (CSR) initiatives
  \item Environmental, social, and governance (ESG) factors
\end{itemize}

\subsection{Quantitative Analysis}

Quantitative analysis involves evaluating a company's financial health through its financial statements. Financial Modelling is the process of creating a mathematical representation of a company's financial performance, often using historical data to forecast future results.

\subsubsection{Interpreting Financial Statements}

\textbf{Types of Financial Statements}
\begin{itemize}
    \item \textbf{Standalone Statements:} Represent the financial position of a single entity e.g., subsidiaries.  So, if a company has multiple businesses, then a standalone statement of the company will offer details of the financial performance of that particular company.
    \item \textbf{Consolidated Statements:} Represent the financial position of the parent company along with its subsidiaries,  associates, or joint ventures. Hence, if you are analyzing the statements of a large organization, then a consolidated statement offers a better picture of the performance of all the companies put together.
\end{itemize}

\textit{Note:} Analyzing consolidated financial statements is preferred for in-depth analysis, but they should be reviewed alongside standalone statements. If a company has many subsidiaries, the standalone is too narrow; consolidated gives the full financial picture, and is generally preferred for evaluating the business comprehensively.

\textbf{Example:} When analyzing a company like Tata Power, it's important to look at both standalone and consolidated financial statements. Standalone statements reflect the financials of Tata Power alone, excluding its subsidiaries, while consolidated statements include all group entities under Tata Power, such as its renewable and distribution arms. Therefore, for a complete view of the company's performance and financial health, consolidated statements are more relevant. It's also incorrect to assume that Tata Power's consolidated statements include Tata Steel, as each is a separate listed entity under the Tata Group.

\subsubsection*{(i) Profit and Loss / Income Statement}

\begin{itemize}
    \item Generally released every quarter (i.e., every 3 months).
    \item Used to analyze the performance of the company during a specific period, typically a financial year.
    \item Reflects the company's operating performance in a given year.
\end{itemize}

\subsection*{Format Overview}

\renewcommand{\arraystretch}{1.5}
\setlength{\tabcolsep}{10pt}

\begin{longtable}{|>{\raggedright\arraybackslash}p{2.5cm}|>{\raggedright\arraybackslash}p{6cm}|>{\raggedright\arraybackslash}p{4.5cm}|}
\hline
\textbf{Component} & \textbf{Description and Formula} & \textbf{What to Look for / Analysis} \\
\hline
\textbf{I. Revenue from Operations / Sales} & Revenue earned from the firm's core activities (e.g. sales of products/services). & Consistent growth indicates a healthy and scalable core business. Compare YoY and QoQ. \\
\hline
\textbf{II. Other Income} & Non-operational income (e.g. interest, dividend, investment gains, foreign exchange gains). & High other income may indicate reliance on non-core earnings - check sustainability. \\
\hline
\textbf{A. Total Revenue} & \( \text{Total Revenue} = \text{I} + \text{II} \) & Used to assess the complete earning picture. Sudden jumps: investigate source. \\
\hline
\textbf{III. Cost of Goods Sold / Operating Expenses} & Direct costs to produce goods/services (e.g. raw materials, labor). & Rising COGS vs stagnant revenue reduces margins - assess pricing power and input costs. \\
\hline
\textbf{B. Gross Profit} & \( \text{Gross Profit} = \text{A} - \text{III} \) & Indicates basic profitability from operations. Monitor gross margin \((\text{Gross Profit} / \text{Revenue})\). \\
\hline
\textbf{IV. Marketing, Advertising, and Promotion Expenses} & Costs to promote the business and drive sales. & High marketing costs may reduce profits but can build brand and drive future growth. \\
\hline
\textbf{V. Sales, General, and Administrative Expenses (SG\&A)} & Indirect operational expenses: salaries, rent, insurance, office costs, etc. Sometimes also include Depreciation and Amortization & Efficiency in managing these overheads improves operating leverage. \\
\hline
\textbf{VI. Other Expenses} & R\&D, foreign exchange loss, impairments, stock-based compensation, utility expenses, CSR expenses, etc. & Volatile or non-recurring items - adjust for normalized earnings. Watch for R\&D investment. \\
\hline
\textbf{C. EBITDA (Earnings Before Interest, Taxes, Depreciation, and Amortization)} & \( \text{EBITDA} = \text{B} - (\text{IV} + \text{V} + \text{VI}) \) & Key measure of core profitability before non-cash and financing impacts. Compare across peers. \\
\hline
\textbf{VII. Depreciation and Amortization} & Non-cash expenses for asset cost allocation. Depreciation: tangible; Amortization: intangible. & High D\&A suggests capital-intensive business. Useful for comparing EBITDA vs EBIT. \\
\hline
\textbf{D. EBIT (Earnings Before Interest and Tax)} & \( \text{EBIT} = \text{C} - \text{VII} \) & Reflects profitability after accounting for non-cash costs. Good proxy for operating performance. \\
\hline
\textbf{E. Operating Profit} & \( \text{Operating Profit} = \text{D} - \text{II} \) & Strips out non-operational income. Reveals true operating health. \\
\hline
\textbf{VIII. Interest Expenses / Finance Costs} & Cost of servicing debt (interest payments). & High interest: assess debt burden. Check interest coverage ratio \( \left( \frac{\text{EBIT}}{\text{Interest}} \right) \). \\
\hline
\textbf{F. PBT (Profit Before Tax)} & \( \text{PBT} = \text{E} - \text{VIII} + \text{II} \) & Profit before taxes. Key for comparing pre-tax margins. Tax holiday or deferred tax can distort comparison. \\
\hline
\textbf{IX. Income Tax / Provision for Tax} & Taxes charged on pre-tax profit. & Effective tax rate \( = \text{Tax} / \text{PBT} \). Analyze tax planning, deferred tax. \\
\hline
\textbf{G. PAT (Profit After Tax)} & \( \text{PAT} = \text{F} - \text{IX} + \text{II} \) & Net earnings after tax. Track net profit margin and YoY PAT growth. \\
\hline
\textbf{H. Net Profit} & \( \text{Net Profit} = \text{G} \) & Final earnings figure. Compare with industry average. Adjust for extraordinary items. \\
\hline
\textbf{I. EPS (Earnings Per Share)} & \( \text{EPS} = \frac{\text{H}}{\text{Total Outstanding Shares}} \) & Key shareholder metric. Use for valuation ratios (P/E). Check for dilution. \\
\hline

\end{longtable}

\subsubsection*{(ii) Balance Sheet}

\begin{itemize}
    \item Generally released on the last day of every fiscal quarter (i.e., every 3 months).
    \item Contains detailed info of the company's assets, liabilities, and shareholders' equity.
    \item \textbf{ASSETS} = \textbf{LIABILITIES} + \textbf{SHAREHOLDERS' EQUITY}
    \item Assets generate revenue $\rightarrow$ revenue minus costs becomes retained earnings $\rightarrow$ retained earnings increase equity. Liabilities are either taken to fund assets or meet short-term obligations. The difference between assets and liabilities is the value attributable to owners (i.e., equity).
    \item Balance sheets should also be compared with those of other businesses in the same industry since different industries have unique approaches to financing.
\end{itemize}


\subsection*{Format Overview}

\renewcommand{\arraystretch}{1.5}
\setlength{\tabcolsep}{6pt}

\begin{longtable}{|>{\raggedright\arraybackslash}p{3cm}|>{\raggedright\arraybackslash}p{7cm}|>{\raggedright\arraybackslash}p{5cm}|}
\hline
\textbf{Component} & \textbf{Description (with formulas)} & \textbf{What to Look for / Analysis} \\
\hline
\multicolumn{3}{|p{15cm}|}{\centering\textbf{EQUITY (Company's actual worth)}} \\
\hline
\textbf{I. Shareholder Equity / Book Value of Company} & Book value of equity contributed by shareholders. It is the paid-up capital from primary issuance like IPO, right issues, etc. & Indicates how much capital has been injected by shareholders; steady increase indicates investor confidence. \\
\hline
\textbf{II. Reserve and Surplus / Retained Earnings} & Cumulative profits made by the company since inception minus dividends paid. & Consistent growth suggests profitability and internal reinvestment capacity. Declines may indicate high dividend payouts or losses. \\
\hline
\textbf{A. Total Shareholder's Fund} & $A = \text{I} + \text{II}$ \newline What would be a company's worth if it goes for sale after clearing all its liabilities. & A strong base capital. Should ideally be increasing over time. Useful for calculating Return on Equity (RoE). \\
\hline
\multicolumn{3}{|p{15cm}|}{\centering\textbf{LIABILITIES (What company owes)}} \\
\hline
\multicolumn{3}{|p{15cm}|}{\centering\textbf{NON CURRENT LIABILITIES (Obligations due beyond 1 year)}} \\
\hline
\textbf{III. Long-term Debt / Borrowings} & Loans and borrowings with a tenure more than 1 year. & Check for rising debt levels. Used for leverage analysis (Debt-to-Equity). High debt signals financial risk. \\
\hline
\textbf{IV. Deferred Tax Liabilities} & Taxes due for current period but payable in future. Non-cash liability - does not affect immediate cash flow but affects future tax payments. & Indicates timing differences between accounting and tax profits. Stable DTLs are normal. \\
\hline
\textbf{V. Long-term Provisions} & Money set aside for employee benefits, gratuity, etc. & Watch for consistency. Sudden increases might signal rising employee obligations. \\
\hline
\textbf{VI. Other Long-term Liabilities} & Includes deposits like security or rent deposits, deferred credits, etc. & Examine footnotes for breakdown. Evaluate credit terms extended to the firm. \\
\hline
\textbf{B. Total Long-term Liabilities} & $B = \text{III} + \text{IV} + \text{V} + \text{VI}$ & Analyze long-term solvency. Should be compared against non-current assets. \\
\hline
\multicolumn{3}{|p{15cm}|}{\centering\textbf{CURRENT LIABILITIES (Obligations to be met within 1 year)}} \\
\hline
\textbf{VII. Short-term Debt / Borrowings} & Debt with tenure less than 1 year. & High values indicate liquidity risk. Crucial for liquidity ratio analysis (e.g., Current Ratio). \\
\hline
\textbf{VIII. Trade Payables} & Amounts due to suppliers for goods and services. & Increasing payables may improve short-term liquidity but could hurt supplier relationships. \\
\hline
\textbf{IX. Other Current Liabilities} & Unclaimed \& Unpaid dividends, overdrafts, current maturities of long-term borrowings. & Understand breakups; higher values may indicate cash crunch. \\
\hline
\textbf{X. Short-term Provisions} & Set aside for short-term obligations like product warranty, employee benefits, gratuity. & High provisions may be prudent or could hint at underlying risks. \\
\hline
\textbf{C. Total Current Liabilities} & $C = \text{VII} + \text{VIII} + \text{IX} + \text{X}$ & Helps assess short-term financial health. Key input in working capital analysis. \\
\hline
\textbf{D. Total Capital and Liabilities} & $D = A + B + C$ & Should match total assets. Forms the liability side of the balance sheet. \\
\hline
\multicolumn{3}{|p{15cm}|}{\centering\textbf{ASSETS (What company owns)}} \\
\hline
\multicolumn{3}{|p{15cm}|}{\centering\textbf{NON CURRENT FIXED ASSETS (assets used in the operation of the business and not intended for resale and can't be easily encashed within 1 year)}} \\
\hline
\textbf{XI. Tangible Assets} & Physical assets like land, building, equipment. & Useful for calculating asset turnover ratio. Helps determine capital intensity. \\
\hline
\textbf{XII. Intangible Assets} & Non-physical assets: patents, brands, goodwill. & Assess sustainability of competitive advantage. Check amortization schedules. \\
\hline
\textbf{XIII. Capital Work-in-Progress} & Costs incurred on fixed assets under construction. & Large CWIP could indicate expansion or delayed projects. \\
\hline
\textbf{XIV. Intangible Assets under Development} & In-progress intangible developments like R\&D, tech platforms. & Examine carefully-may offer future benefits or signify speculative investment. \\
\hline
\textbf{E. Total Fixed Assets} & $E = \text{XI} + \text{XII} + \text{XIII} + \text{XIV}$ & Forms backbone of productive capacity. Compare with long-term debt. \\
\hline
\multicolumn{3}{|p{15cm}|}{\centering\textbf{NON CURRENT NON FIXED ASSETS (assets not actively used in operations but held for strategic or financial purposes and can't be easily encashed within 1 year)}} \\
\hline
\textbf{XV. Non-current Investments} & Investments held for long term (e.g., in subsidiaries, associates). & Strategic holdings. Review for returns, synergy, and stability. \\
\hline
\textbf{XVI. Long-term Loans and Advances} & Loans provided to others with maturity beyond one year. & Must check creditworthiness of borrowers. Could be non-core. \\
\hline
\textbf{XVII. Other Non-current Assets} & Assets not classified above (e.g., deferred expenses, deposits). & Inspect for materiality and recurring nature. \\
\hline
\textbf{F. Total Non-current Assets} & $F = E + \text{XV} + \text{XVI} + \text{XVII}$ & Total long-term resource allocation. Compare with long-term liabilities. \\
\hline
\multicolumn{3}{|p{15cm}|}{\centering\textbf{CURRENT ASSETS (assets expected to be converted into cash, sold, or consumed within one year or the business's operating cycle)}} \\
\hline
\textbf{XVIII. Current Investments} & Investments liquid within 1 year (e.g., mutual funds, treasury bills). & High values support liquidity. May earn short-term income. \\
\hline
\textbf{XIX. Inventories} & Items held for resale or production. & Key metric for manufacturing/retail. Watch inventory turnover ratio. \\
\hline
\textbf{XX. Cash and Cash Equivalents} & Includes bank balance, cash on hand, liquid deposits. & Key liquidity buffer. Compare with short-term liabilities. \\
\hline
\textbf{XXI. Short-term Loans and Advances} & Loans repayable within a year. & Examine credit quality and collection trends. \\
\hline
\textbf{XXII. Other Current Assets} & e.g., trade receivables, accrued income. & Analyze receivables aging. High values could mean cash is tied up. \\
\hline
\textbf{G. Total Current Assets} & $G = \text{XVIII} + \text{XIX} + \text{XX} + \text{XXI} + \text{XXII}$ & Key to working capital analysis. Crucial for short-term financial decisions. \\
\hline
\textbf{H. Total Assets} & $H = F + G$ & Should match $D$ (Total Capital and Liabilities). Use to calculate RoA, asset turnover. \\
\hline

\end{longtable}

\subsubsection*{(iii) Cash Flow Statement}
\begin{itemize}
    \item Generally released every quarter (i.e., every 3 months).
    \item Provides insights into cash inflows and outflows from operating, investing, and financing activities.
    \item Helps assess liquidity, solvency, and financial flexibility.
    \item One of the most important financial statements, as it shows how well a company generates cash to pay its debt obligations and fund its operating expenses. Also it can not be manipulated easily.
\end{itemize}


\subsection*{Format Overview}

\renewcommand{\arraystretch}{1.5}
\setlength{\tabcolsep}{6pt}

\begin{longtable}{|>{\raggedright\arraybackslash}p{3cm}|>{\raggedright\arraybackslash}p{7cm}|>{\raggedright\arraybackslash}p{5cm}|}
\hline
\textbf{Component} & \textbf{Description (with formulas)} & \textbf{What to Look for / Analysis} \\
\hline
\multicolumn{3}{|c|}{\textbf{Cash Flow from Operating Activities}} \\
\hline
\textbf{I. Net Profit Before Tax} & Profit before considering taxes and extraordinary items. & Starting point for operating cash flow. Should reconcile with the income statement. \\
\hline
\textbf{II. Depreciation and Amortization} & Non-cash expense for allocation of cost of assets over time. & Added back since it reduces profit but not cash. Higher D\&A in asset-heavy firms. \\
\hline
\textbf{III. Adjustments} & Includes non-operating gains/losses (e.g., investment gains), deferred taxes, and stock-based compensation. & Crucial for adjusting accrual-based profit to reflect cash reality. Should be monitored for quality of earnings. \\
\hline
\textbf{IV. Working Capital Changes} & $\Delta WC = \Delta \text{Current Assets} - \Delta \text{Current Liabilities}$ & Indicates how much cash is tied up in operations. Large increases in working capital reduce cash flow. \\
\hline
\textbf{V. Interest Paid} & Cash outflows due to interest payments on borrowings. & Deducted under operating cash flows (as per Indian GAAP). Indicates cost of debt. \\
\hline
\textbf{VI. Tax Paid} & Actual tax outflow during the year. & High tax outflow reduces cash; differs from tax expense in P\&L due to timing differences. \\
\hline
\textbf{A. Operating Cash Flow} & $A = \text{I} + \text{II} \pm \text{III} \pm \text{IV} - \text{V} - \text{VI}$ & Most important segment for cash sustainability. Should be positive and growing over time. \\
\hline
\multicolumn{3}{|c|}{\textbf{Cash Flow from Investing Activities}} \\
\hline
\textbf{VII. Net Capital Expenditures} & Cash used to acquire or upgrade fixed assets (e.g., plant, machinery). & Usually negative. Shows long-term reinvestment. Excessive capex without returns is a red flag. \\
\hline
\textbf{VIII. Net Investments} & Investments in financial securities or divestments from subsidiaries. & Evaluate whether investments are productive or speculative. Watch for cash inflow from asset sales. \\
\hline
\textbf{B. Investing Cash Flow} & $B = \pm \text{VII} \pm \text{VIII}$ & Ideally negative due to capital reinvestment. Sustained positive values may mean asset stripping or low growth. \\
\hline
\multicolumn{3}{|c|}{\textbf{Cash Flow from Financing Activities}} \\
\hline
\textbf{IX. Dividends Paid} & Cash paid to shareholders as dividend. & Reflects return to equity holders. High payout may reduce reinvestment ability. \\
\hline
\textbf{X. Sale or Purchase of Company Stock} & Inflows from issuing equity or outflows from buybacks. & Dilution (issue) vs confidence (buybacks). Indicates capital raising or capital return. \\
\hline
\textbf{XI. Net Borrowings} & Cash flows from issuing or repaying loans. & High borrowing inflow could mean cash need; repayment indicates deleveraging. \\
\hline
\textbf{C. Financing Cash Flow} & $C = -\text{IX} \pm \text{X} \pm \text{XI}$ & Direction of flow indicates whether company is raising or returning capital. \\
\hline
\textbf{Net Cash Flow} & $= A + B + C$ & Should reconcile with change in cash on balance sheet. Positive net cash indicates cash accumulation. \\
\hline

\end{longtable}


\subsubsection*{$\rightarrow$ How to Analyze Financial Statements?}

\textbf{Types of Analysis}
\begin{description}
    \item[1. Vertical Analysis (Common Size Statements):]
    Each line item is expressed as a percentage of a base figure, usually total revenue (in income statements) or total assets (in balance sheets).

    \[
    \text{Vertical Analysis} = \left( \frac{\text{Statement Line Item}}{\text{Total Base Figure}} \right) \times 100
    \]

    \textit{Example:} If Revenue = \rupee1,00,000 and COGS = \rupee40,000, then:
    \[
    \frac{40,000}{1,00,000} \times 100 = 40\% \quad \text{(COGS as \% of Revenue)}
    \]

    Useful for comparing companies of different sizes and analyzing internal cost structure.\\

    \item[2. Horizontal Analysis:]
    Measures the change in a line item over time to detect trends or growth.

    \[
    \text{Horizontal Analysis} = \left( \frac{\text{Current Year Amount} - \text{Base Year Amount}}{\text{Base Year Amount}} \right) \times 100
    \]

    \textit{Example:} If Revenue in 2022 = \rupee1,00,000 and in 2023 = \rupee1,20,000, then:
    \[
    \frac{1,20,000 - 1,00,000}{1,00,000} \times 100 = 20\% \quad \text{(Revenue growth)}
    \]

    Useful for trend analysis and assessing year-on-year or quarter-on-quarter performance.
\end{description}

\textbf{What to Look for while analyzing Financial Statements?}

\begin{enumerate}
    \item \textbf{Operating Profit} is the profit a company generates from its core business functions, where interest and tax deduction is excluded
    \begin{align}
    \text{Operating Profit} &= \text{Net Sales} - (\text{COGS} + \text{Administrative} \\
    &\quad + \text{Overheads} + \text{Depreciation} + \text{Amortization})
    \end{align}
    
    \item \textbf{Net Profit (Bottom Line)} = Total Revenue - Total Expenses
    
    \item \textbf{EBITDA} = Net Profit + Interest + Tax + Depreciation \& Amortization
    \begin{itemize}
      \item Better measure for Capital Intensive Industries (Manufacturing, Oil \& Gas, Telecom...)
    \end{itemize}
    \textbf{EBIT} = Net Profit + Interest + Tax
    \begin{itemize}
        \item Better measure for Service Industries (Technology, Consulting)
    \end{itemize}
    
    \item \textbf{Other Income}: Company can try to inflate its net Income by increasing its income from other activities.
    
    \item \textbf{Depreciation and Amortization}: Sometimes company changes its accounting policies for calculation of Depreciation \& Amortization to inflate its net income.
    
    Two methods to calculate:
    \begin{enumerate}[label=(\roman*)]
        \item Straight Line method
        \item Accelerated method
    \end{enumerate}
    
    \item \textbf{Debt/Borrowing}: Lower the debt is better but if company is utilizing the debt the right then its fine. Increasing debt and high interest rates can be very dangerous for the company.
    
    \item \textbf{Reserves and Surplus}: It can be good for a company to face financial crisis but too much is not good.
    
    \item \textbf{Trade Payables} is the amount payable by a business for goods purchased or services availed as a part of their business operations. \\
    \textbf{Trade Receivables} is just the opposite i.e., amount receivable.
    
    \item \textbf{Cash Flow from Operations} summarizes all cash inflows and outflows from an organization's operational activities.
    \begin{align}
    \text{CFO} &= \text{PAT} + \text{D\&A} \pm \text{Changes in Current Assets} \\
    &\quad \pm \text{Changes in Current Liabilities}\\
    &= \text{Total Revenue} - \text{Operating Expenses}
    \end{align}
    
    \text{Note:} Positive Cashflow from operations is better. 10-20\% less CFO than PAT is ok but else is not so good.
    
    \item \textbf{Cash Flow from Investment} - These represent depend on cash inflow and outflow used from investing activities of an organisation.
    \begin{align}
    \text{CFO} &= \text{Purchase/Sale of Long term Assets} \\
    &\quad + \text{Other Businesses} + \text{Marketable Securities}
    \end{align}

    \text{Note:} Negative Cashflow from investments is better.
    
    \item \textbf{Cash Flow from Financing Activities} - represents all cash inflow and outflow from financing activities
    \begin{align}
    \text{CFF} &= \text{Issue/Repurchase (Equity + Debt)} \\
    &\quad - \text{Dividend Payments and other items}
    \end{align}

    \text{Note:} Negative Cashflow from financing is better.
    
    \item \textbf{Net Cash Flow} is total inflow of cash. (+ve, better)
    Question any increase/decrease/change in important line items and why it changes from previous year/quarter.
\end{enumerate}


\subsubsection*{$\rightarrow$ How to Read Annual Reports of a Company?}

Annual Report is a financial document published annually to showcase a company's business operations and financial condition for the concerned fiscal year.

\textbf{Key Sections to Focus On}

\begin{enumerate}
    \item \textbf{Starting Pages} $\rightarrow$ Vision, branding (should not be too much)
    \item \textbf{Product's Suite}
    \item \textbf{Key Performance Indicators}
    \item \textbf{Compliances}
    \item \textbf{Chairman's Letter}: Its objective is to state the reasons regarding how a company addressed its failures and related issues. Provides in-depth analysis of the overall business performance. Indicates challenges and risks faced by company. Conclude with financial plans and business direction for upcoming years.
    \item \textbf{MD's Message}
    \item \textbf{Risk Mitigation}
    \item \textbf{Board of Directors}
    \item \textbf{Business Profile}: Revenue generated, details about products/services manufactured/provided, subsidiaries owned, capital market in which it competed, risk factors, company policy, new products/services, etc.
    \item \textbf{Financial Statements}
    \item \textbf{Notes to Financial Statements}
    \item \textbf{Management Discussion and Analysis}: See management is answering their questions, avoiding or correcting it?
    \item \textbf{Industry Insight}
    \item \textbf{Share Holding Pattern}
    \item \textbf{Debt Analysis}
    \item \textbf{Board Meetings}
    \item \textbf{Salary of Management}
    \item \textbf{Related Party Transactions}
    \item \textbf{Corporate Governance}
\end{enumerate}

\subsubsection*{$\rightarrow$ How to Read and Understand Quarterly Results?}

Quarterly Reports are summary of unaudited financial statements issued by company every quarter (3 months i.e., 4 times a year - June, Sept, Dec, March)

\textbf{Quarterly Reports analysis}

\begin{itemize}
    \item Helps to access the performance of the company invested in or to make informed decisions
    \item Consist of Operational Highlights, Press release, FAQ sheet, Quarterly Presentations, Shareholding Pattern, Earning Call Transcript, Management Commentary
    \item It can show the seasonality or trend of the company and industry. You should find reasons for the same
    \item Compare the Quarterly Results to previous one and to the previous year's same quarter
    \item If you find some manipulations in the results you should see the inside story and treat it like a red flag $\triangleright$
\end{itemize}



\subsection*{$\rightarrow$ Red Flag Analysis of a Company?}

Identifying red flags is crucial in avoiding fundamentally weak or potentially fraudulent companies. Watch out for the following warning signs:

\begin{enumerate}
    \item \textbf{Declining Revenues} \\
    Consistent fall in top-line growth indicates weak demand or operational issues.
    
    \item \textbf{Rising Debt-to-Equity Ratio} \\
    Excessive leverage raises financial risk and reduces solvency.
    
    \item \textbf{Volatile Cash Flows} \\
    Unpredictable cash inflows/outflows reflect operational instability.
    
    \item \textbf{Stock Price Fluctuation} \\
    Extreme or unexplained volatility in share price.
    
    \item \textbf{Pending Legal Cases} \\
    Ongoing litigation could result in financial or reputational damage.
    
    \item \textbf{Falling Promoter Holding} \\
    Continuous decrease signals lack of confidence in the business.
    
    \item \textbf{High Promoter Pledging} \\
    Rising pledging of shares is a risk to ownership stability.
    
    \item \textbf{High NPAs (especially for Banks)} \\
    Indicates poor lending practices and financial deterioration.
    
    \item \textbf{Delayed Reports} \\
    Late filing of annual or quarterly reports may suggest governance issues.
    
    \item \textbf{Over-Optimistic Financials} \\
    Results that appear too good to be true often are.
    
    \item \textbf{Auditor Red Flags} \\
    Any mention of "misstatements" or exceptions in audit reports. Read the auditor's report ("Summary of misstatements") carefully.
    
    \item \textbf{Unusual Accounting Practices} \\
    Frequent changes in accounting methods or policies.
    
    \item \textbf{Changes in Financial Reporting Style} \\
    May indicate manipulation or hiding of actual performance.
    
    \item \textbf{Sudden Management Changes} \\
    Frequent turnover raises stability concerns.
    
    \item \textbf{Last-Minute Bulk Sales} \\
    Sales booked mostly in the last few days of a period are suspicious.
    
    \item \textbf{Asset Revaluation Surprises} \\
    Sudden large changes in fixed/intangible asset value.
    
    \item \textbf{Complex Transactions} \\
    Unnecessarily complicated deals may hide risks or fraud.
    
    \item \textbf{Performance-Linked Bonuses} \\
    Incentives tied to manipulated metrics can misalign interests.
    
    \item \textbf{Strange behaviour of key Ratios} \\
    This analysis should cover the following key areas:
    \begin{itemize}
        \item Gross Profit Margin \(\uparrow\) but Sales \(\downarrow\)
        \item Return on Equity (ROE) \(\downarrow\) but Price-to-Book (P/B) Ratio \(\uparrow\)
    \end{itemize}
    
    \item \textbf{Rising Debt or Inventory} \\
    Build-up in debt or unsold goods implies weak demand or inefficiency.
\end{enumerate}

\textit{Note: Multiple red flags together strengthen the case for caution or avoidance. Always investigate the underlying causes.}

\subsubsection{Ratio Analysis}

Ratio analysis is a quantitative method of evaluating a company's financial performance by comparing various financial metrics. It helps investors and analysts assess profitability, liquidity, solvency, and efficiency.

\subsubsection*{(i) Solvency / Financial Risk Ratios}

Used to measure an entrise's ability to meet its long term debt obligations.

\renewcommand{\arraystretch}{1.5}
\begin{longtable}{|>{\raggedright\arraybackslash}m{4.5cm}|>{\raggedright\arraybackslash}m{5cm}|>{\raggedright\arraybackslash}m{6cm}|}
\hline
\textbf{Ratio Name} & \textbf{Formula} & \textbf{Analysis \& Interpretation} \\
\hline
Debt to Equity Ratio & 
$\frac{\text{Total Debt}}{\text{Total Equity}}$ &
\begin{itemize}
\item Measures financial leverage
\item How much business is funded by debt compared to equity distribution
\item Ideal: $<1$ 
\item Compare with industry peers (e.g., acceptable $>2$ for capital-intensive industries like Power, Banks) 
\end{itemize} \\
\hline

Debt Ratio & 
$\frac{\text{Total Debt}}{\text{Total Assets}}$ &
\begin{itemize}
\item Shows \%age of assets financed by debt
\item Safe threshold: $<0.5$
\item $>0.7$ indicates aggressive leverage and high risk
\item Critical for cyclical industries 
\item Generally hight debt ratio in industries like Power, Banks but with stable cash flows, yet should not be $>0.75$
\end{itemize} \\
\hline

Debt to Capital Ratio & 
$\frac{\text{Total Debt}}{\text{Debt + Shareholders' Equity}}$ &
\begin{itemize}
\item shows the debt company has compared to its equity
\item Similar analysis to Debt Ratio
\item Ideal: $<0.5$
\end{itemize} \\
\hline

Interest Coverage Ratio (ICR) & 
$\frac{\text{EBIT}}{\text{Interest Expense}}$ &
\begin{itemize}
\item Debt servicing capacity
\item How many times a company can cover its current interest payments with its available earnings
\item Safe: $>3$ for most industries, $<1.5$ indicates distress risk
\item For utilities like Power companies: $>2$  and $<3$ is recommended
\end{itemize} \\
\hline

Debt Service Coverage Ratio (DSCR) & 
$\frac{\text{Operating Profit}}{\text{Total Debt Service}}$ &
\begin{itemize}
\item Cash flow adequacy for debt payments
\item Healthy: $1.5<$ DSR $<2$
\item Debt Service = interest + principal repayments 
\end{itemize} \\
\hline

Financial Leverage Ratio & 
$\frac{\text{Total Assets}}{\text{Shareholders' Equity}}$ &
\begin{itemize}
\item How much the assets belongs to shareholders than debt creditors
\item Higher ratio = More debt reliance (Lower better)
\end{itemize} \\
\hline

Proprietary Ratio (Equity Ratio) & 
$\frac{\text{Shareholders' Funds}}{\text{Total Assets}}$ &
\begin{itemize}
\item To what extent Shareholders' funds have been invested in the assets of the business
\item Ideal: $>0.5$; $<0.3$ indicates high creditor stake
\end{itemize} \\
\hline
\end{longtable}

\subsubsection*{(ii) Liquidity Ratios}

Used to determine a debtor's ability to pay off current debt obligations without raising external capital.


\renewcommand{\arraystretch}{1.5}
\begin{longtable}{|>{\raggedright\arraybackslash}m{4.5cm}|>{\raggedright\arraybackslash}m{5cm}|>{\raggedright\arraybackslash}m{6cm}|}
\hline
\textbf{Ratio Name} & \textbf{Formula} & \textbf{Analysis \& Interpretation} \\
\hline
Current Ratio & 
$\frac{\text{Current Assets}}{\text{Current Liabilities}}$ &
\begin{itemize}
\item Short-term debt paying ability
\item Ideal: 1.33 - 3 
\item $>3$: Inefficient asset utilization, funds are either idle or locked up in inventories and receivables
\item $<1$: Potential liquidity crisis
\end{itemize} \\
\hline

Quick Ratio (Acid Test) & 
$\frac{\text{Current Assets} - \text{Inventory}}{\text{Current Liabilities}}$ &
\begin{itemize}
\item Strict liquidity measurement
\item Ideal: $1<$ Cash Ratio $<2.5$
\item $>2.5$ suggests improper utilization of short term funds
\item Critical for industries with slow-moving inventory
\end{itemize} \\
\hline

Cash Ratio & 
$\frac{\text{Cash + Cash Equivalents}}{\text{Current Liabilities}}$ &
\begin{itemize}
\item How quick and to what extent, company can repay its current dues with help of its readily available assets
\item Ideal: 0.5 - 1
\item $>1$: Excessive idle cash
\item $<0.2$: High default risk
\end{itemize} \\
\hline
\end{longtable}

\subsubsection*{(iii) Efficiency Ratios}

Measures company's ability to use its assets and manage its liabilities effectively in the current period or in short term.

\renewcommand{\arraystretch}{1.5}
\begin{longtable}{|>{\raggedright\arraybackslash}m{4.5cm}|>{\raggedright\arraybackslash}m{5cm}|>{\raggedright\arraybackslash}m{6cm}|}
\hline
\textbf{Ratio Name} & \textbf{Formula} & \textbf{Analysis \& Interpretation} \\
\hline
Inventory Turnover Ratio & 
\begin{itemize}
  \item $\frac{\text{Cost of Goods Sold}}{\text{Average Inventory}}$
  \item \text{Inventory Turnover Days} = $\frac{365}{\text{Inventory Turnover Ratio}}$
\end{itemize} &
\begin{itemize}
\item Measures inventory management efficiency
\item Higher = Better (industry dependent), Manufacturing: 6-8x ideal, Retail: 12-15x desirable
\item Inventory Turnover Days tells on average how many days the company's inventory is held until it is sold. Generally Lower better, but much low means company is not fulfilling its demand.
\end{itemize} \\
\hline

Receivable Turnover Ratio & 
\begin{itemize}
  \item $\frac{\text{Net Credit Sales}}{\text{Average Account Receivables}}$
  \item \text{Account Receivable Days} = $\frac{365}{\text{Receivable Turnover Ratio}}$
\end{itemize} &
\begin{itemize}
\item How quick company collects bills from its customers, Higher = Faster conversion
\item Net Credit Sales = Sales on credit - Sales returns - Sales Allowances
\item Account Receivable Days is how many days does the company takes to collect cash from its credit sales from the customers
\end{itemize} \\
\hline

Payable Turnover Ratio & 
\begin{itemize}
  \item $\frac{\text{COGS}}{\text{Average Payables}}$
  \item \text{Payable Days} = $\frac{365}{\text{Payable Turnover Ratio}}$
\end{itemize} &
\begin{itemize}
\item Payment discipline measurement
\item Lower = Better credit terms, but high canmean the company has favourable credit terms
\item Payable Days is the average payment period
\end{itemize} \\
\hline

Total Asset Turnover Ratio & 
$\frac{\text{Net Sales}}{\text{Average Total Assets}}$ &
\begin{itemize}
\item Ability of company to generate sales from its assets
\item Higher = Better capital use
\item 0.5-2x = Common range, For Capital-intensives: 0.3-0.8x normal
\end{itemize} \\
\hline

Working Capital Turnover & 
\begin{itemize}
  \item $\frac{\text{Net Sales}}{\text{Average Working Capital}}$
  \item Working Capital = Current (Assets - Liabilities)
  \item Working Capital Cycle = (Inventory + Receivable - Payable) Days
\end{itemize} &
\begin{itemize}
\item Reflects the amount of operating profit needed to maintain a given level of sales
\item Higher = Better (5-10x ideal)
\item Working Capital is the amount of cash business can safely spending
\item Working Capital / Cash Conversion Cycle is the time it takes to convert total net working capital into cash. Ideally should be negative
\end{itemize} \\
\hline
\end{longtable}

\subsubsection*{(iv) Financial / Profitability Risk Ratios}

\textbf{Growth Ratios}

Indicates how fast the business is growing, Higher better.

\renewcommand{\arraystretch}{1.5}
\begin{longtable}{|>{\raggedright\arraybackslash}m{4.5cm}|>{\raggedright\arraybackslash}m{5cm}|>{\raggedright\arraybackslash}m{6cm}|}
\hline
\textbf{Ratio Name} & \textbf{Formula} \\
\hline

Operating Profit Margin & 
$\frac{\Delta \text{Operating Profit}}{\text{Net Sales}} \times 100$ \\
\hline

Gross Profit Margin & 
$\frac{\Delta \text{Gross Profit}}{\text{Net Sales}} \times 100$ \\
\hline

EPS Growth & 
$\frac{\Delta \text{EPS}}{\text{Prior EPS}} \times 100$ \\
\hline
\end{longtable}

Similarly, we can calculate EBITDA margin, PAT margin, etc.

We can also calculate their Compounded Annual Growth Rate (CAGR) or simple percentage change as follows, where \(n\) is the number of periods (e.g., years):

\[
\text{CAGR} = \left( \frac{\text{End Value}}{\text{Start Value}} \right)^{\frac{1}{n}} - 1
\]
\[
\text{\% change} = \frac{\text{Final value} - \text{Initial Value}}{\text{Initial Value}} \times 100
\]

\textbf{Return Ratios}

\renewcommand{\arraystretch}{1.5}
\begin{longtable}{|>{\raggedright\arraybackslash}m{4.5cm}|>{\raggedright\arraybackslash}m{5cm}|>{\raggedright\arraybackslash}m{6cm}|}
\hline
Return on Investment (ROI) & 
$\frac{\text{Net Return}}{\text{Cost}} \times 100$ &
\begin{itemize}
\item It is just a basic return measurement for any investment you make.
\end{itemize} \\
\hline

Return on Invested Capital (ROIC) & 
$\frac{\text{EBIT}(1-\text{Tax})}{\text{Debt} + \text{Equity}}$ &
\begin{itemize}
\item How well a company is using its money to generate returns
\item ROIC > WACC = Value creation
\item If the company uses debt to finance investments, calculate WACC to determine if ROIC exceeds the cost of capital (including the interest rate on debt), indicating value creation.
\end{itemize} \\
\hline

Return on Equity (ROE) & 
$\frac{\text{PAT}}{\text{Total Equity}} \times 100$ &
\begin{itemize}
\item How well is the firm in generating returns on the investment received from its shareholders
\item Higher Better
\item High ROE with high debt = Risk, Companies can inflate ROE by increasing it's debt.
\end{itemize} \\
\hline

Return on Capital Employed (ROCE) & 
$\frac{\text{EBIT}}{\text{Capital Employed}}$ &
\begin{itemize}
\item Overall fficiency to generate profits from capital
\item Capital Employed = Equity + Debt (Long Term + Short Term)
\item Higher Better, ROCE > WACC
\end{itemize} \\
\hline

Return on Assets (ROA) & 
$\frac{\text{EBIT or PAT}}{\text{Total Assets}}$ &
\begin{itemize}
\item Measures how efficiently a company uses its assets to generate earnings
\item Higher ROA is better
\item ROA can be increased by:
    \begin{itemize}
        \item Increasing EBIT (better pricing, economies of scale, cost efficiency)
        \item Decreasing assets (reducing receivables, better inventory turnover, higher asset utilization, divesting low-margin businesses, leasing instead of buying)
    \end{itemize}
\end{itemize} \\
\hline

Weighted Average Cost of Capital (WACC) & 
$\frac{\text{Equity}}{\text{Total Capital}} \times \text{Cost of Equity OR Expected rate of Return} + \frac{\text{Debt}}{\text{Total Capital}} \times \text{Interest Rate} \times (1 - \text{Tax Rate})$ &
\begin{itemize}
\item Represents a firm's blended cost of capital (equity + debt)
\item Required return for a capital budgeting project
\item Lower is better for investors; higher implies a company must generate more returns
\end{itemize} \\
\hline
\end{longtable}

\textbf{Cash Flow Ratios}

\renewcommand{\arraystretch}{1.5}
\begin{longtable}{|>{\raggedright\arraybackslash}m{4.5cm}|>{\raggedright\arraybackslash}m{5cm}|>{\raggedright\arraybackslash}m{6cm}|}
\hline
Cash Flow Per Share & 
$\frac{\text{Operating CF} - \text{Pref. Div}}{\text{Outstanding Shares}}$ &
\begin{itemize}
\item Real cash generation, it is the net cash a company produces on a per share basis
\item Compare with EPS
\item CFPS > EPS = Quality earnings
\end{itemize} \\
\hline

Free Cash Flow Per Share (FCFPS) & 
$\frac{\text{Free Cash Flow}}{\text{Total Outstanding Shares}} = \frac{\text{Operating CF} - \text{CAPEX}}{\text{Outstanding Shares}}$ &
\begin{itemize}
\item Indicates a company's ability to pay debt, dividends, buy back stock, and grow the business
\item Higher FCFPS is better
\item Increasing FCF implies potential increase in share price
\item FCF = cash left after operating expenses and asset-related capital expenditure
\item Capital Expenditure (CAPEX) = funds used to acquire or upgrade physical assets like plant, property, and equipment
\end{itemize} \\
\hline

Free Cash Flow to Equity (FCFE) & 
$\text{PAT} + \text{Depr. \& Amor.} - \text{CAPEX} \pm \text{Change in current assets} \pm \text{Change in current liabilities}$ &
\begin{itemize}
\item Cash available to equity holders after meeting all debt obligations
\end{itemize} \\
\hline

Free Cash Flow to Company (FCFF) & 
$\text{EBIT}(1-\text{Tax}) + \text{D\&A} - \text{CAPEX} \pm \text{Change in current assets} \pm \text{Change in current liabilities}$ &
\begin{itemize}
\item Cash flow available to equity and debt holders in the company
\end{itemize} \\
\hline

Dividend Yield & 
$\frac{\text{Annual Div/Share}}{\text{Stock Price}} \times 100$ &
\begin{itemize}
\item measures how much a firm pays out in dividents each year relative to stock price
\item 2-4\% = Healthy; >6\% = Sustainability check
\end{itemize} \\
\hline
\end{longtable}

\subsubsection*{$\rightarrow$ Dupont Analysis of ROE}

The DuPont analysis is a financial framework that decomposes Return on Equity (ROE) into three key components to provide insights into a company's profitability, operational efficiency, and financial structure. By breaking down ROE, it helps investors and analysts understand whether high returns are driven by strong profits, efficient asset use, or significant debt, making it a valuable tool for fundamental analysis.

\textbf{Formula}

The DuPont analysis expresses ROE as the product of three ratios:

\[
\text{ROE} = \left( \frac{\text{Net Income}}{\text{Sales}} \right) \times \left( \frac{\text{Sales}}{\text{Assets}} \right) \times \left( \frac{\text{Assets}}{\text{Shareholders' Equity}} \right)
\]

Where:
\begin{itemize}
    \item \textbf{Profit Margin} (\(\frac{\text{Net Income}}{\text{Sales}}\)): Measures how much profit a company generates per dollar of sales, reflecting its ability to control costs and pricing power.
    \item \textbf{Asset Turnover} (\(\frac{\text{Sales}}{\text{Assets}}\)): Indicates how efficiently a company uses its assets to generate sales, showing operational efficiency.
    \item \textbf{Financial Leverage} (\(\frac{\text{Assets}}{\text{Shareholders' Equity}}\)): Represents the degree to which a company uses debt to finance its assets, amplifying returns but increasing risk.
\end{itemize}

This decomposition allows for a detailed analysis of the factors contributing to a company's ROE, facilitating comparisons across firms and industries.

\textbf{Components of DuPont Analysis}

\begin{itemize}
    \item \textbf{Profit Margin}: A high profit margin indicates strong profitability, often due to premium pricing, cost efficiency, or a favorable product mix. For example, technology companies like Apple often have high margins due to brand strength.
    \item \textbf{Asset Turnover}: A higher ratio suggests efficient use of assets to generate revenue. Retail companies typically have high asset turnover due to rapid inventory turnover, while capital-intensive firms may have lower ratios.
    \item \textbf{Financial Leverage}: A high leverage ratio indicates reliance on debt, which can boost ROE but also increases financial risk. Companies with stable cash flows can sustain higher leverage safely.
\end{itemize}

\textbf{Steps in DuPont Analysis}

\begin{enumerate}
    \item \textbf{Collect Financial Data}: Obtain Net Income, Sales (Revenue), Total Assets, and Shareholders' Equity from the company's financial statements (income statement and balance sheet).
    \item \textbf{Calculate Profit Margin}: Divide Net Income by Sales to determine profitability per dollar of revenue.
    \item \textbf{Calculate Asset Turnover}: Divide Sales by Total Assets to assess asset efficiency.
    \item \textbf{Calculate Financial Leverage}: Divide Total Assets by Shareholders' Equity to measure debt usage.
    \item \textbf{Compute ROE}: Multiply the three ratios to obtain ROE and verify by calculating ROE directly (\(\frac{\text{Net Income}}{\text{Shareholders' Equity}}\)).
    \item \textbf{Analyze Results}: Interpret each component to understand the drivers of ROE and compare with industry peers.
\end{enumerate}

\textbf{Example: Apple Inc. (Fiscal Year Ended September 28, 2024)}

To illustrate the DuPont analysis, consider Apple Inc.'s financial data for the fiscal year ended September 28, 2024, sourced from its annual report:

\begin{table}[h]
\centering
\begin{tabular}{|l|r|}
\hline
\textbf{Metric} & \textbf{Value (in millions)} \\
\hline
Net Income & \$93,736 \\
Total Revenue (Sales) & \$391,035 \\
Total Assets & \$364,980 \\
Shareholders' Equity & \$56,950 \\
\hline
\end{tabular}
\caption{Financial Data for Apple Inc., Fiscal Year 2024}
\end{table}

Using these figures, we calculate each component of the DuPont formula:

1. \textbf{Profit Margin}:
   \[
   \text{Profit Margin} = \frac{\text{Net Income}}{\text{Sales}} = \frac{93,736}{391,035} \approx 0.2397 \text{ or } 23.97\%
   \]
   Apple generates approximately 24 cents in profit for every dollar of sales, reflecting strong profitability due to its premium pricing and brand strength.

2. \textbf{Asset Turnover}:
   \[
   \text{Asset Turnover} = \frac{\text{Sales}}{\text{Assets}} = \frac{391,035}{364,980} \approx 1.0714
   \]
   Apple generates \$1.07 in sales per dollar of assets, indicating moderate efficiency in asset utilization.

3. \textbf{Financial Leverage}:
   \[
   \text{Financial Leverage} = \frac{\text{Assets}}{\text{Shareholders' Equity}} = \frac{364,980}{56,950} \approx 6.4081
   \]
   For every dollar of equity, Apple has \$6.41 in assets, suggesting significant use of debt or liabilities, though mitigated by substantial cash reserves.

Combining these, the ROE is:
\[
\text{ROE} = 0.2397 \times 1.0714 \times 6.4081 \approx 1.6458 \text{ or } 164.58\%
\]

To verify, calculate ROE directly:
\[
\text{ROE} = \frac{\text{Net Income}}{\text{Shareholders' Equity}} = \frac{93,736}{56,950} \approx 1.6458 \text{ or } 164.58\%
\]

The calculations align, confirming the accuracy of the DuPont analysis.

\textbf{Interpretation}

The DuPont analysis reveals that Apple's exceptionally high ROE of 164.58% is driven by:
- \textbf{High Profit Margin (23.97\%)}: Apple's ability to command premium prices for its products (e.g., iPhones, MacBooks) and maintain cost efficiency results in strong profitability.
- \textbf{Moderate Asset Turnover (1.0714)}: Apple uses its assets efficiently to generate sales, though not as aggressively as retail or manufacturing firms with higher turnover ratios.
- \textbf{High Financial Leverage (6.4081)}: Apple's low shareholders' equity relative to assets, partly due to share buybacks, amplifies ROE. However, its large cash reserves (e.g., \$29.97 billion in cash as of September 2024) reduce the risk associated with high leverage.

This analysis highlights Apple's strong profitability and strategic use of leverage, tempered by moderate asset efficiency. Investors can compare these components to industry peers to assess relative performance.


\subsubsection*{(v) Valuation Ratios}

Helps to develop a sense of how the stock price is valued by market participants


\renewcommand{\arraystretch}{1.5}
\begin{longtable}{|>{\raggedright\arraybackslash}m{4.5cm}|>{\raggedright\arraybackslash}m{5cm}|>{\raggedright\arraybackslash}m{6cm}|}
\hline
\textbf{Ratio Name} & \textbf{Formula} & \textbf{Analysis \& Interpretation} \\
\hline

Earning Per Share (EPS) & 
$\frac{\text{Net Income - Preffered Dividents}}{\text{Outstanding Shares}}$ &
\begin{itemize}
\item To assess a company's profitability on a per-share basis
\item A rising EPS over time signals improving profitability, a positive sign for investors.
\end{itemize} \\
\hline

Price-to-Earnings (P/E) & 
$\frac{\text{Market Price}}{\text{EPS}}$ &
\begin{itemize}
\item How much the market is willing to pay for every rupee of company's current earnings
\item Sector averages vary widely, Compare with peers, industry average and historical P/E of the company.
\item Low P/E = stock might be undervalued, low or negative growth expectation
\item High P/E = Growth expectations, stock might be overvalued
\end{itemize} \\
\hline

Price-to-Book (P/B) & 
\begin{itemize}
  \item $\frac{\text{Market Price}}{\text{Book Value/Share}}$
  \item Book value = Total assets - Intangible assets - Liabilities
  \item Book value per share = $\frac{\text{Book value}}{\text{Outstanding Shares}}$
\end{itemize} &
\begin{itemize}
\item How much shareholders are paying for a company's net assets
\item Book value per share indicates a firm's net asset value on per share basis
\item High P/B ratio indicate market is ready to pay premium on book value because of future growth of company or value of intangible assets
\item More useful for asset heavy industries (and not for technology and service industries)
\end{itemize} \\
\hline

PEG Ratio & 
$\frac{\text{P/E}}{\text{EPS Growth Rate}}$ &
\begin{itemize}
\item Denotes rupee amount a shareholder needs to pay for a stock of a company for earning \rupee 1 of its income
\item <1 = Undervalued
\item >1.5 = Overvalued
\item not a goog metric for asset heavy businesses
\end{itemize} \\
\hline

Price-to-Cash Flow & 
$\frac{\text{Market Price}}{\text{Operating Cash Flow/Share}}$ &
\begin{itemize}
\item Rupee value an investor is willing to pay for the cash flow generated
\item Compare with P/E ratio, Lower better
\item if ther is a gap between net profit and operating cash flow, the company should be investing money in expansion and not in renewal
\end{itemize} \\
\hline

EV/EBITDA & 
$\frac{\text{Enterprise Value}}{\text{EBITDA}}$ &
\begin{itemize}
\item Capital structure neutral, used to compare the entire value of a business with the amount of EBITDA it earns on an annual basis
\item 8-12x = Fair value range
\item <6x = Acquisition target
\end{itemize} \\
\hline

Price-to-Sales (P/S) & 
$\frac{\text{Market Cap}}{\text{Revenue}}$ &
\begin{itemize}
\item Shows how much themarket values every rupee of compan's sales, Higher better
\item Compare with peers
\item Best ratio to watch if company is at temporary losses.
\end{itemize} \\
\hline
\end{longtable}

\subsubsection*{(vi) Ratios For Banks}

\renewcommand{\arraystretch}{1.5}
\begin{longtable}{|>{\raggedright\arraybackslash}m{4.5cm}|>{\raggedright\arraybackslash}m{5cm}|>{\raggedright\arraybackslash}m{6cm}|}
\hline
\textbf{Ratio} & \textbf{Formula} & \textbf{Analysis} \\
\hline

CASA (Current Account Saving Account) Ratio & 
$\frac{\text{CASA Deposits}}{\text{Total Deposits}}$ &
\begin{itemize}
\item Indicates how much of total bank deposits are in current and savings accounts
\item Higher is better (lower cost of funds), But not always, balance needed
\item Typically, banks offer no interest on current a/c, 2--3\% on savings a/c
\end{itemize} \\
\hline

Credit to Deposits \% (LTD) & 
$\frac{\text{Bank's Total Loan Disbursed}}{\text{Total Deposits Received}}$ &
\begin{itemize}
\item Indicates how much of deposits are lent out
\item Too low: bank not earning much
\item Too high: liquidity risk in unforeseen events
\item Ideal between 80\% - 90\%
\end{itemize} \\
\hline

Interest Expended to Total Funds \% & 
$\frac{\text{Interest Expense}}{\text{Total Deposits}}$ &
\begin{itemize}
\item Shows what portion of income is spent on interest for borrowed money
\item Lower is better
\item $<$10\% implies sufficient liquidity
\item $>$10\% is generally not good
\end{itemize} \\
\hline

Interest Income to Total Funds \% & 
$\frac{\text{Net Interest Income}}{\text{Total Deposits in Bank}}$ &
\begin{itemize}
\item Measures bank's earning from interest relative to total funds
\item Net Interest Income = Interest Earned - Interest Paid
\item Higher is better
\end{itemize} \\
\hline

Advances Growth \% & 
$\frac{\text{Last Year Adv.} - \text{This Year Adv.}}{\text{Last Year Advances}} \times 100$ &
\begin{itemize}
\item Shows how much bank advances have grown year-on-year
\item Advances = loans given to individuals/businesses to meet working capital needs
\item Higher is better
\end{itemize} \\
\hline

Non-Performing Assets (NPA) & 
$\text{GNPA\%} = \frac{\text{Gross NPA}}{\text{Total Advances}} \times 100$ \\
& 
$\text{NNPA\%} = \frac{\text{Net NPA}}{\text{Total Advances}} \times 100$ &
\begin{itemize}
\item Loans overdue by more than 90 days are NPAs
\item Lower is better, shows better loan book quality
\item Gross NPA is the total value of advances that have become NPA in a particular quarter/financial year
\item Net NPA = GNPA - Provisions (actual bad loans after accounting provisions)
\item Provisions are the funds set aside by Company/Bank as assets to pay for anticipated future losses/taxes.
\end{itemize} \\
\hline

Cost to Income Ratio \% & 
$\frac{\text{Operating Expenses}}{\text{Operating Income}}$ &
\begin{itemize}
\item Measures efficiency: cost incurred to earn income
\item Lower is better
\end{itemize} \\
\hline

Operating Cost to Assets Ratio \% & 
$\frac{\text{Operating Expenses}}{\text{Average Assets}}$ &
\begin{itemize}
\item Measures operating efficiency of bank relative to asset base
\item Lower is better - more profitable and efficiently managed
\end{itemize} \\
\hline

Cost of Liabilities \% & 
$\frac{\text{Total Interest on Borrowings and Deposits}}{\text{Total Borrowings and Deposits}}$ &
\begin{itemize}
\item Indicates rate paid on current debt
\item Lower is better — improves margins and profitability
\end{itemize} \\
\hline

Capital Adequacy Ratio (CAR) & 
$\frac{\text{Total Capital}}{\text{Risk Weighted Assets}}$ &
\begin{itemize}
\item Shows bank's ability to absorb potential losses
\item Higher is better - protects depositor money and ensures stability
\end{itemize} \\
\hline

Net Interest Margin (NIM) Ratio & 
$\frac{\text{Interest Earned} - \text{Interest Paid}}{\text{Assets that Bank Holds}}$ &
\begin{itemize}
\item Measures profitability from lending/investment activities
\item Must be positive; higher is better
\end{itemize} \\
\hline

Yield on Advances \% & 
$\frac{\text{Total Interest Earned}}{\text{Average Loans \& Advances}}$ &
\begin{itemize}
\item Measures return bank earns on its loans
\item Higher is better
\end{itemize} \\
\hline
\end{longtable}

\subsubsection*{(vii) Market Meter Ratios}


\renewcommand{\arraystretch}{1.5}
\begin{longtable}{|>{\raggedright\arraybackslash}m{4.5cm}|>{\raggedright\arraybackslash}m{5cm}|>{\raggedright\arraybackslash}m{6cm}|}
\hline
\textbf{Ratio Name} & \textbf{Formula} & \textbf{Analysis \& Interpretation} \\
\hline
Market-Cap to GDP (Buffett Indicator) & 
$\frac{\text{Total Market Cap}}{\text{GDP of country}} \times 100$ &
\begin{itemize}
\item Broad market valuation
\item For India:
\begin{itemize}
\item $>75\%$: Overvalued (Sell signal)
\item $<75\%$: Undervalued (Buy signal)
\end{itemize}
\item US Historical Avg: 100-120\%
\end{itemize} \\
\hline

Debt to GDP Ratio & 
$\frac{\text{Total Government Debt}}{\text{GDP of country}} \times 100$ &
\begin{itemize}
\item Sovereign risk measure (country's ability to repay its debt)
\item By Debt here we mean, the fiscal deficit (the difference between the total income of government and its total expenditure)
\item India Target: $<64\%$
\item $>64\%$: can lead to Market crash
\item Developed Nations: $80-120\%$ common
\end{itemize} \\
\hline

Nifty/Sensex P/E Ratio & 
$\frac{\text{Total Market Cap of Nifty 50}}{\text{Total Earnings of Nifty 50}}$ &
\begin{itemize}
\item Market valuation gauge
\item Compare with 10-Yr Average
\item India Historical Range: 18-28 P/E
\item $>25$: Caution warranted
\end{itemize} \\
\hline
\end{longtable}

\subsubsection{Valuation Analysis}

\subsubsection*{Intrinsic Value}
\textbf{Intrinsic Value} is the anticipated or calculated value of a stock, currency, or product based on fundamentals rather than market price.

\textbf{Key Points to Remember}
\begin{enumerate}[label=\arabic*.]
    \item It's all about future cash flows.
    \item You must discount the cash flows to the present day.
    \item Understand the business model clearly.
    \item It gives a rough figure not an exact number.
    \item Always begin with assumptions.
\end{enumerate}

\subsubsection*{(i) Method 1: Asset-Based Valuation}
\textbf{Formula:}
\[
\text{Value per share}
\]
\[
= \frac{\text{Market Capitalization of Company} + \sum \left( \text{Percentage Holding in company} \times \text{Market Cap of company} \right)}{\text{Total Number of Equity Shares in Parent / Holding Company}}
\]
\[
= \frac{\text{Net Assets} - \text{Preference Shares Capital} - \text{Debt}}{\text{Number of Equity Shares}}
\]

\textbf{Use Case:} Primarily for holding/parent companies with multiple assets.  
\textbf{Note:} Suggests buying when there's a \textit{discount rate} (e.g., 25\%).

\subsubsection*{(ii) Method 2: Market Based Valuation}
\begin{itemize}
    \item \textbf{Expected rate of Dividents:}
    \[
    \frac{\text{Profit available for Dividents}}{\text{Paid-up Equity Capital}} \times 100
    \]
    \item \textbf{Expected rate of Earnings:}
    \[
    \frac{\text{Profit After Tax}}{\text{Paid-up Equity Capital}} \times 100
    \]
    \item \textbf{Value per Share:}
    \[
    \frac{\text{Expected rate of Earnings or Dividents}}{\text{Normal Rate of Return}} \times {\text{Equity paid up Capital}}
    \]
\end{itemize}

\subsubsection*{(iii) Method 3: Comparable Company Analysis (Relative Valuation)}
\begin{itemize}
    \item Aims to derive a stock's theoretical price using the price multiples of similar companies
    \item It involves the use of valuation ratios like P/E ratio, P/B Ratio, EV/EBITDA, EV/Sales, etc.
    \item Compare company with its historical and Industrial Averages and with its peers.
\end{itemize}

\subsubsection*{(iv) Method 4: Dividend Discount Model (DDM)}
\textbf{Use Case:} Companies that pay stable dividends.

\textbf{Formula (Gordon Growth Model):}
\[
\text{Fair Value} = \frac{D_1}{r - g}
\]
Where:
\begin{itemize}
    \item \( D_1 \) = Expected dividend next year
    \item \( r \) = Cost of equity
    \item \( g \) = Dividend growth rate
\end{itemize}

\subsubsection*{(v) Method 5: Discounted Cash Flow (DCF)}

The Discounted Cash Flow (DCF) method is a valuation technique used to estimate the intrinsic value of a company, project, or asset based on its expected future cash flows. It incorporates the time value of money, recognizing that a dollar today is worth more than a dollar in the future due to its potential earning capacity. The DCF method is widely used in investment finance, corporate financial management, and for valuing companies, particularly for long-term investment decisions \cite{InvestopediaDCF}.

\textbf{Formula}

The DCF value is calculated as the sum of the present values of future cash flows over a forecast period, often combined with a terminal value to account for cash flows beyond that period:

\[
\text{DCF} = \sum_{t=1}^{n} \frac{\text{CF}_t}{(1 + r)^t} + \frac{\text{TV}}{(1 + r)^n}
\]

Where:
\begin{itemize}
    \item \(\text{CF}_t\): Cash flow in year \(t\),
    \item \(r\): Discount rate,
    \item \(n\): Number of forecast periods,
    \item \(\text{TV}\): Terminal value, representing cash flows beyond the forecast period.
\end{itemize}

\textbf{Components of DCF}

\begin{itemize}
    \item \textbf{Cash Flows (\(\text{CF}_t\))}: These are typically free cash flows to the firm (FCFF), which represent cash available to all investors (debt and equity holders) after operating expenses, taxes, and capital expenditures. Alternatively, free cash flows to equity (FCFE) can be used for equity valuation, reflecting cash available to shareholders after debt payments. FCFF is preferred for enterprise valuation, as it captures the total value of the company.
    \item \textbf{Discount Rate (\(r\))}: The rate used to discount future cash flows to their present value, reflecting the investment's risk. For FCFF, the weighted average cost of capital (WACC) is commonly used, calculated as the weighted average of the cost of equity and cost of debt. For FCFE, the cost of equity is used. The WACC accounts for the company's capital structure and risk profile.
    \item \textbf{Terminal Value (\(\text{TV}\))}: Since companies are assumed to generate cash flows indefinitely, the terminal value estimates the present value of all cash flows beyond the forecast period. It is often calculated using the perpetuity growth model:
    \[
    \text{TV} = \frac{\text{FCFF}_{n+1}}{r - g}
    \]
    Where \(\text{FCFF}_{n+1}\) is the cash flow in the year after the forecast period, and \(g\) is the perpetual growth rate (typically 2–3\%, reflecting long-term economic growth).
\end{itemize}

\textbf{Steps in DCF Valuation}

\begin{enumerate}
    \item \textbf{Forecast Free Cash Flows}: Estimate FCFF or FCFE for a period (e.g., 5–10 years) based on financial statements, revenue growth, operating margins, and capital expenditures.
    \item \textbf{Determine the Discount Rate}: Calculate the WACC using the cost of equity (e.g., via the Capital Asset Pricing Model) and cost of debt, weighted by the company's capital structure.
    \item \textbf{Calculate Present Value of Cash Flows}: Discount each year's cash flow to the present using the formula \(\frac{\text{CF}_t}{(1 + r)^t}\).
    \item \textbf{Estimate Terminal Value}: Use the perpetuity growth model or an exit multiple (e.g., based on EBITDA) to estimate cash flows beyond the forecast period.
    \item \textbf{Discount Terminal Value}: Discount the terminal value to the present using \(\frac{\text{TV}}{(1 + r)^n}\).
    \item \textbf{Sum Present Values}: Add the discounted cash flows and terminal value to get the enterprise value.
    \item \textbf{Adjust for Equity Value}: Subtract net debt (total debt minus cash) from the enterprise value to obtain the equity value. Divide by the number of shares outstanding to get the per-share value.
\end{enumerate}

\textbf{Example}

Consider a company with the following projected FCFF over five years, a WACC of 5\%, and an initial investment of \$11 million (e.g., for acquiring the company):

\begin{table}[h]
\centering
\begin{tabular}{|c|c|c|}
\hline
\textbf{Year} & \textbf{Cash Flow} & \textbf{Discounted Cash Flow} \\
\hline
1 & \$1,000,000 & \$952,381 \\
2 & \$1,000,000 & \$907,029 \\
3 & \$4,000,000 & \$3,455,350 \\
4 & \$4,000,000 & \$3,290,810 \\
5 & \$6,000,000 & \$4,701,157 \\
\hline
\end{tabular}
\caption{Discounted Cash Flows for the Company}
\end{table}

The present value of each cash flow is calculated as:

\[
\text{PV of CF}_t = \frac{\text{CF}_t}{(1 + r)^t}
\]

For each year:
\begin{itemize}
    \item Year 1: \(\frac{1,000,000}{(1 + 0.05)^1} = 952,381\)
    \item Year 2: \(\frac{1,000,000}{(1 + 0.05)^2} = 907,029\)
    \item Year 3: \(\frac{4,000,000}{(1 + 0.05)^3} = 3,455,350\)
    \item Year 4: \(\frac{4,000,000}{(1 + 0.05)^4} = 3,290,810\)
    \item Year 5: \(\frac{6,000,000}{(1 + 0.05)^5} = 4,701,157\)
\end{itemize}

Summing these gives the total DCF for the forecast period:

\[
\text{Total DCF} = 952,381 + 907,029 + 3,455,350 + 3,290,810 + 4,701,157 = 13,306,727
\]

To account for cash flows beyond year 5, assume a perpetual growth rate (\(g\)) of 2\%. The FCFF in year 6 is estimated as \(\text{FCFF}_6 = \text{FCFF}_5 \times (1 + g) = 6,000,000 \times 1.02 = 6,120,000\). The terminal value is:

\[
\text{TV} = \frac{6,120,000}{0.05 - 0.02} = \frac{6,120,000}{0.03} = 204,000,000
\]

Discount the terminal value to the present (year 0):

\[
\text{PV of TV} = \frac{204,000,000}{(1 + 0.05)^5} = \frac{204,000,000}{1.2762815625} \approx 159,863,013
\]

The enterprise value is:

\[
\text{Enterprise Value} = 13,306,727 + 159,863,013 = 173,169,740
\]

If the company has \$2 million in net debt, the equity value is:

\[
\text{Equity Value} = 173,169,740 - 2,000,000 = 171,169,740
\]

Assuming 1 million shares outstanding, the fair price per share is:

\[
\text{Price per Share} = \frac{171,169,740}{1,000,000} \approx 171.17
\]

The Net Present Value (NPV) for the investment decision is:

\[
\text{NPV} = 173,169,740 - 11,000,000 = 162,169,740
\]

Since the NPV is positive, the investment is worthwhile.

\section{Sentiment Analysis}

\subsection{Key Factors Affecting Market Sentiment}

\begin{enumerate}
    \item \textbf{Change in Open Interest (OI)}\\
    Open interest refers to the total number of outstanding derivative contracts (like futures and options) that haven't been settled.
    \begin{itemize}
        \item Rising OI with price increase suggests new long positions (bullish).
        \item Rising OI with price decrease suggests new short positions (bearish).
        \item Falling OI with price movement suggests position closure (profit booking).
    \end{itemize}

    \item \textbf{Delivery Percentage}\\
    Indicates the portion of traded stocks actually delivered to investors (not just intraday trading).
    \begin{itemize}
        \item High delivery percentage implies strong investor confidence (bullish sentiment).
        \item Low delivery indicates speculative or intraday activity (less reliable sentiment).
    \end{itemize}

    \item \textbf{FII/DII Buying and Selling}\\
    Foreign Institutional Investors (FII) and Domestic Institutional Investors (DII) activity significantly influences market trends.
    \begin{itemize}
        \item Net buying by FIIs generally boosts sentiment.
        \item Net selling indicates bearishness or profit booking.
    \end{itemize}

    \item \textbf{News}\\
    Market sentiment is heavily influenced by financial, political, and global news.
    \begin{itemize}
        \item Positive news (mergers, results, economic growth) leads to bullish trends.
        \item Negative news (scandals, poor earnings, wars) leads to bearish trends.
    \end{itemize}

    \item \textbf{Moving Average}\\
    A moving average smooths out price data to identify trends over time.
    \begin{itemize}
        \item Crossing above a key moving average (e.g., 50-day or 200-day) indicates upward momentum.
        \item Crossing below it indicates downward momentum.
    \end{itemize}

    \item \textbf{Volatility Index (VIX)}\\
    The India VIX measures the market's expectation of volatility over the next 30 days for Nifty.
    \begin{itemize}
        \item Low VIX implies low fear and stable sentiment.
        \item High VIX implies high fear and uncertainty.
        \item Generally, \textbf{VIX} and \textbf{Nifty} move inversely:
        \begin{itemize}
            \item If VIX $\downarrow$ $\rightarrow$ Nifty $\uparrow$ (bullish)
            \item If VIX $\uparrow$ $\rightarrow$ Nifty $\downarrow$ (bearish)
        \end{itemize}
        \item If VIX $> 22$, it indicates high volatility best to avoid fresh investments, especially in derivatives (F\&O).
        \item \[\text{VIX} / 2\sqrt{2} \, \% \, \rightarrow \text{gives the percentage increase or decrease for next 30 days}\]
    \end{itemize}

    \item \textbf{Institutional Activity}:
    \[
    \text{FII/DII Ratio} = \frac{\text{Foreign Institutional Investment}}{\text{Domestic Institutional Investment}}
    \]
    \begin{itemize}
        \item $>1.2$ = Bullish foreign outlook
        \item $<0.8$ = Domestic confidence
    \end{itemize}

    \item \textbf{Sentiment Analysis Framework}
    \begin{center}
    \begin{tabular}{|l|l|l|}
    \hline
    \textbf{Tool} & \textbf{Bullish Signal} & \textbf{Bearish Signal} \\ \hline
    Put/Call Ratio & $<0.7$ & $>1.3$ \\ \hline
    Short Interest & $<5\%$ & $>20\%$ \\ \hline
    News Sentiment & $70\%+$ Positive & $70\%+$ Negative \\ \hline
    Google Trends & Rising "Buy" searches & Rising "Sell" searches \\ \hline
    \end{tabular}
    \end{center}

    \item \textbf{Research Reports from Reputed Brokerages}\\
    Reports and recommendations from reputed financial institutions influence investor confidence.
    \begin{itemize}
        \item Positive ratings/upgrades can create bullish sentiment.
        \item Negative ratings/downgrades can trigger sell-offs.
    \end{itemize}

    \item \textbf{Government Policies}\\
    Policy decisions related to taxation, infrastructure, or regulations have a direct impact on sectors and overall sentiment.
    \begin{itemize}
        \item Investor-friendly policies = bullish trend.
        \item Restrictive or uncertain policies = bearish trend.
    \end{itemize}

    \item \textbf{Ignore Online Forums}\\
    Online forums often spread unverified information or hype.
    \begin{itemize}
        \item Avoid basing trades on unsubstantiated social media discussions.
        \item Rely on data-driven or institutional analysis instead.
    \end{itemize}

    \item \textbf{Investor Interest using Trading Volume}\\
    Trading volume helps assess investor interest and market strength.
    \begin{itemize}
        \item High volume with price rise = strong bullish sentiment.
        \item High volume with price drop = strong bearish pressure.
    \end{itemize}

    \item \textbf{Stock Price}\\
    The price action of a stock reflects real-time sentiment.
    \begin{itemize}
        \item Consistent price gains show confidence.
        \item Declines suggest caution or negativity.
    \end{itemize}
\end{enumerate}

\subsection{Price Manipulation Techniques}

Stock price manipulation involves artificially inflating or deflating the price of a security to deceive or influence investors.

\begin{enumerate}
    \item \textbf{Increasing/Decreasing Demand or Supply}\\
    Manipulators can create false buying or selling pressure to move the stock price. For instance, placing large buy orders can increase demand and raise prices temporarily.

    \item \textbf{Pledging of Shares by Promoters}\\
    A very high or increasing percentage of promoter-pledged shares may be a red flag. It may lead to price volatility if lenders liquidate shares during price drops.

    \item \textbf{Low Daily Traded Volume and Delivery Percentage}\\
    Stocks with low liquidity are easier to manipulate. A small trade can lead to significant price movement, making them vulnerable to artificial price setting.

    \item \textbf{Volume Changes Without Price Changes}\\
    Anomalies such as a surge in trading volume without a corresponding price move can indicate accumulation or distribution by manipulators.

    \item \textbf{Manipulation by High Net-Worth or Retail Investors}\\
    Even retail investors or HNIs, acting in groups or syndicates, can manipulate illiquid stocks through coordinated buying/selling.

    \item \textbf{Fundamentally Weak Stocks Are Easier to Manipulate}\\
    Stocks lacking strong financials or public scrutiny are more prone to manipulation as fewer investors track or question unusual activity.
\end{enumerate}

\subsection{Operator Activity}

\textbf{Operators} are individuals or entities who artificially drive market prices using tactics such as pump and dump schemes. 

\begin{itemize}
    \item They often form syndicates involving brokers, insiders, or speculators.
    \item The aim is to manipulate prices for personal gain while appearing as normal market activity.
    \item They create buying pressure without price movement to accumulate stock, causing others to enter when price starts to rise.
    \item Once retailers enter, they exit by dumping shares, causing the price to fall.
\end{itemize}

\textbf{Typical Activity Pattern}

\begin{itemize}
    \item \textbf{Entry:} $\uparrow$ Volume $\rightarrow$ No Price Change $\rightarrow$ Price likely to rise $\uparrow$
    \item \textbf{Exit:} $\uparrow\uparrow$ Volume $\rightarrow$ No Price Change $\rightarrow$ Price likely to fall $\downarrow$
\end{itemize}

\textbf{Operator Activity Patterns}
\begin{center}
\begin{tabular}{|l|l|l|}
\hline
\textbf{Phase} & \textbf{Volume Pattern} & \textbf{Price Action} \\ \hline
Accumulation & Low volume & Sideways movement \\ \hline
Markup & Gradual increase & Steady upward trend \\ \hline
Distribution & High volume & Volatile swings \\ \hline
Dump & Massive volume & Sharp decline \\ \hline
\end{tabular}
\end{center}

\section{Investment Psychology and Portfolio Management}

\subsection{Behavioral Risks}
  \begin{itemize}
  \item \textbf{Ego:} Almost everyone tends to believe they are better than average in various aspects of life, particularly in interpersonal skills. This is known as the \textbf{``better-than-average effect''}, a well-documented cognitive bias in psychology. In the context of investing, this overconfidence can be particularly dangerous.
  \textbf{Reality Check:} The majority of investors cannot consistently outperform the market. According to efficient market hypothesis (EMH), beating the market requires above-average skill, but by definition, not everyone can be above average. There must be a distribution of returns with some investors performing at or below the average.

  \item \textbf{Coservatism:} refers to the tendency of investors to underreact to new information, remaining anchored to prior beliefs even when the evidence suggests otherwise. This results in delayed or suboptimal decision-making, particularly in fast-changing markets.
  \textbf{Confirmation Bias} further amplifies this issue. Investors often search for, interpret, and recall information in ways that confirm their existing beliefs or hypotheses. For instance, they may focus on the positive aspects of a stock they already own (if long) and ignore warning signs, or vice versa if they are short.
  The psychological root of these biases lies in our discomfort with cognitive dissonance the mental discomfort we experience when presented with contradictory information. The ability to accept new information, self-reflect, and admit being wrong is rare but essential for successful investing.

  \item \textbf{Attention:} Humans are naturally drawn to vivid narratives, especially during uncertain times. \textbf{Attention bias} refers to the tendency to give disproportionate weight to emotionally charged or widely publicized events, while ignoring statistical or probabilistic facts.\\
  \textbf{Example:} During market crashes, financial news media often dramatizes events with headlines like “Markets in Turmoil” or “Economic Apocalypse”. In such times, it's difficult to remember that corrections and bear markets are statistically common, and part of a healthy financial ecosystem.
  \textbf{Behavioral Insight:} This bias is related to the concept of \textbf{availability heuristic} where people assess the probability of events based on how easily examples come to mind. Stories and visuals dominate rational thinking.

  \item \textbf{Emotions:}
  Strong emotions particularly fear and greed are among the most disruptive forces in investing. Emotional decision-making often overrides rational judgment and leads investors to break their own rules.
  \textbf{Example:} An investor who normally sticks to a disciplined asset allocation strategy might panic-sell during a market dip, violating their long-term investment plan.
\end{itemize}

\subsection{Behavioral Biases}

\begin{itemize}
  \item \textbf{Forecasting Errors}
  People overweight recent experiences while making forecasts. This leads to overly optimistic or pessimistic valuations.
  \textbf{Insight:} Investors often extrapolate recent trends into the future, ignoring regression to the mean.

  \item \textbf{Mental Accounting}
  Investors mentally separate their money into different accounts, each with a different investment style. For example, safe investments for a child's education and risky bets elsewhere.
  \textbf{Problem:} This violates the principle of fungibility of money and can lead to suboptimal portfolio management.

  \item \textbf{Regret Avoidance (Realisation Trap)}
  Investors feel more regret when unconventional decisions go wrong. They may stick to "mainstream" investments to avoid blame.
  \textbf{Example:} Losing money in a penny stock is more regretful than losing in a blue-chip stock, even if the amount is the same.

  \item \textbf{Loss Aversion Bias}
  Losses hurt more than gains bring joy. Investors hold onto underperforming stocks hoping to break even, while selling winners too early.
  \textbf{Related Concept:} Average Price Theory averaging down to reduce losses, often leading to poor decision-making.

  \item \textbf{Sunk Cost Fallacy}
  Investors continue to invest in a losing asset to justify past investments.
  \textbf{Example:} Paying a hefty coaching fee and continuing despite poor quality, just because it's already paid.

  \item \textbf{Decision Paralysis}
  Too many choices lead to indecision. Investors may delay decisions to avoid regret.
  \textbf{Note:} Not making a decision is itself a decision. Paralysis can be more costly than making a bad choice.

  \item \textbf{Mental Heuristics}
  The brain takes shortcuts (heuristics) to process information. Investors simplify complex problems.
  \textbf{Example:} Buying a stock because a few factors seem favorable, ignoring deeper analysis (like buying a stock just because it announces Dividents).

  \item \textbf{False Consensus Effect:} People think that most of the people are thinking the same as they do, but that's not the case. There comes a lot of investment opportunities that come in the eyes of few. Also ways of investment, strategies and choices \& hobbies are different for most of the people. Similarly, it's not important that what you are investing in, is invested by the most and vice versa.
  \textbf{Ambiguity Effect:} There are times when we don't know exact result of our decisions, just like investing. So you take views from a lot of people. We should remember that investment strategies are personal and majority opinion is not always right.

  \item \textbf{Herd Mentality and Contrarian Investing}
  Investors follow the crowd, especially under uncertainty.
  \textbf{Contrarian Wisdom:} "Be fearful when others are greedy, and greedy when others are fearful."
  \textbf{Example:} Contrarian investors buy in bear markets when valuations are low and others are panicking.
\end{itemize}

\subsection{Portfolio Management}

Portfolio is a collection of financial investments like Real Estate, stocks, bonds, commodities, cash and cash equivalents, ETFs, mutual funds, etc., and foreign investments.

\textbf{Types of Risks}
\begin{itemize}
\item Uncontrollable Risks:
    \begin{itemize}
    \item Market/Systematic Risk: Uptrend / Downtrend, Macroeconomics
    \item Regulatory Changes (restrictions and policies)
    \end{itemize}
    
\item Controllable Risks:
    \begin{itemize}
    \item Business/Industry Specific (Business cycles, innovations, government regulations)
    \item Portfolio Concentration and Optimization
    \end{itemize}
\end{itemize}

\subsection*{(i) Asset Allocation}
Asset Allocation means matching your risk-return profile, goals, and time horizon with your portfolio's assets.\\

\textbf{Steps:}
\begin{enumerate}
    \item First establish what is your \textbf{risk} and what you expect \textbf{returns} from investments.
    \item Understand the \textbf{goals} you want to conquer and the \textbf{time horizon} of your investment.
    \item Match your risk with your goals.
\end{enumerate}

\noindent\textbf{Note:} Always have some (20--25\%) cash to average in future or better invest.

\subsection*{(ii) Diversification}
Diversification means investing across different asset classes to maximize your returns and minimize your risks.

\begin{quote}
    ``Don't put all your eggs in one basket''
\end{quote}

Diversify your risk-assets properly between high and low risk assets. Even if you are someone who can take risks, it is better to spread your risk and invest across assets that are not related to each other.

\textbf{Optimal / Proper Diversification}
To own a number of individual investments large enough to nearly eliminate unsystematic risk but small enough to concentrate on the best opportunities because different industries, sectors, and assets don't move up and down at the same time or at the same rate.

\textbf{Statistics:}
The average standard deviation / risk / volatility level of a single stock portfolio is 49.2\%, while increasing the number of stocks in the average well-balanced portfolio could reduce it to a max. of 19.2\% (at around 1000 stocks).

\begin{itemize}
    \item A portfolio of \textbf{20 stocks} reduces risk to about 20\%.
    \item The other 980 stocks will reduce it by only 0.8\%.
\end{itemize}

\noindent\textbf{Hence you should not \underline{over diversify}.}

Also, don't invest more than 25\% of your portfolio (in stocks) in a sector.

\subsection*{(iii) Risk Management}
Risk Management involves the identification of threats to an outlay and analysing them through standard mathematical approaches or other means, then eventually deriving measures to mitigate the same.

\textbf{Measures:}
\begin{enumerate}
    \item \textbf{Beta or Alpha} of your portfolio
    \item \textbf{Value at Risk (VaR)} tells you ''How much could I lose in the worst case scenario?'' There are certain calculators where you can calculate it easily.
\end{enumerate}

\subsection*{(iv) Rebalancing}
Rebalancing involves periodically buying/selling assets in a portfolio to maintain an original or desired level of asset allocation to risk (weighted average risk). It also involves \textbf{hedging}. Market and Asset values change all the time. They are affected by a lot of known and unknown factor both domestic and international.\\
\textbf{When to consider Portfolio Rebalancing:}
\begin{itemize}
  \item Portfolio directing away from investor's preference
  \item Change in your circumstances
  \item Poor performance of investment
  \item Global event that shook markets
  \item Better Opportunities
\end{itemize}

\subsection*{$\rightarrow$ When to Sell Stocks}

\textbf{Sell Triggers}
\begin{enumerate}
\item Mistake in Buying Decision
\item Stock Hits Target Price
\item Portfolio Rebalancing Needed
\item Personal Financial Requirements
\item Business Fundamentals Deteriorate:
    \begin{itemize}
    \item Falling Sales/Profits
    \item Rising Debt
    \item Management Issues
    \end{itemize}
    
\item Better Investment Opportunities
\item Sector Hype Ends
\item Corporate Governance Issues
\item Dividends Reduced/Eliminated
\item Lower Trading Volumes
\item Price Rises Too Fast Without Reason
\end{enumerate}

\section{Other Investing and Stock Market Concepts}

\subsection{Market Participants}

\begin{itemize}
    \item \textbf{Promoters:} Individuals or organizations that help raise funds for a business. They may receive shares or a commission for their efforts.
    
    \item \textbf{Retail Investors:} Also called individual investors. These are non-professionals who invest in stocks, mutual funds, ETFs, etc.
    
    \item \textbf{Foreign Institutional Investors (FIIs):} Investment entities registered in a different country that invest in financial markets such as equity, debt, and mutual funds.

    \item \textbf{Foreign Direct Investment (FDI):} Investments made by a company in another country (e.g., setting up a plant or acquiring business assets).

    \item \textbf{Domestic Institutional Investors (DIIs):} Local institutions such as banks, insurance companies, and mutual funds that invest in domestic markets.

    \item \textit{Note:} Generally, FII (investments) $\propto \dfrac{1}{\text{DII (investments)}}$ though both may sometimes invest or divest simultaneously.

    \item \textbf{Stock Exchange:} Platform where companies raise capital by listing shares for public investment. \\
    \textit{Examples:} NSE (National Stock Exchange), BSE (Bombay Stock Exchange).

    \item \textbf{Index:} A statistical measure to track the performance of a group of securities. \\
    \textit{Examples:} Sensex, Nifty 50, Dow Jones.

    \item \textbf{SEBI (Securities and Exchange Board of India):} Regulates the Indian capital market.

    \item \textbf{Stock Broker:} A licensed intermediary who executes buy/sell orders and earns via brokerage and maintenance fees.

    \item \textbf{Market Capitalization Categories:}
    \begin{itemize}
        \item Large Cap: 1st to 100th companies.
        \item Mid Cap: 101st to 250th companies.
        \item Small Cap: 251st onwards.
    \end{itemize}

    \item \textbf{Mutual Funds:} Investment vehicles that pool money from multiple investors to invest in stocks, bonds, and short-term debt.
\end{itemize}

\subsection{Types of Affiliate Companies}

\begin{itemize}
    \item \textbf{Subsidiary:} A company controlled (atleast 50\% ownership) by a parent company to gain synergies (e.g., tax benefits, asset control, diversified risk).

    \item \textbf{Associate:} A company in which a parent owns 20-50\% stake, allowing participation in policy and strategic decisions.

    \item \textbf{Joint Venture (JV):} A contractual business agreement between two or more entities sharing joint control and pooling resources for the purpose of accomplising a specific task, a new project or any other business activity. \textit{Example Ratios:} 50:50, 60:40, 70:30

    \item \textbf{Affiliates:} Companies related through partial ownership (e.g., less than 50\%) or third-party holdings.

    \item \textbf{Mergers:} Two separate companies combine to form a new entity.

    \item \textbf{Acquisitions:} One company takes over another and assumes control.
\end{itemize}

\subsection{Corporate Actions}

\begin{itemize}
    \item \textbf{Dividend} is a reward that publicly listed companies extend to their shareholders from the company's net profit. Companies may decide to retain their accumulated profits to reinvest in the business or reserve it for future use or distribute it to shareholders as dividends. 

    Announcement about dividends makes a significant increase or decrease in the company's stock value. When a company declares a dividend, its share price undergoes a significant increase and starts to decline by a similar proportion once the date of dividend eligibility expires. 

    \textit{Important dates:} Announcement date, Ex-dividend date, Record date, and Payment date.
    
    \item \textbf{Bonus Shares} are free shares that the stockholders of the company receive against shares held by them. Bonus shares are issued out of the reserves in shareholders' equity funds. If the ratio is 1:1, the stockholder receives 1 share for every 1 share held. 

    In bonus issues, the stock price falls in the same proportion as the bonus issue, so we get no real profits. It is issued as an alternative to dividends or to make the shares affordable. 
    
    \textit{Important dates are similar to dividend declaration.}
    
    \item \textbf{Stock Split} is when the number of shares gets multiplied. Technically no new share is issued. If a company announces a stock split in the ratio 1:2, it means for every 1 share held, it will be split into 2 shares.

    Face value of the shares also reduces in the same ratio and so the stock price thus overall no profit. It increases the share liquidity and keeps the share affordable to more shareholders.

    \item \textbf{Reverse Split} is just the reverse of stock split.

    \item \textbf{Rights Issue} is when the company offers new shares to all the existing stockholders in the ratio of their holding in the company at a discount rate to encourage them. This is a primary issue in which the money paid accrues to the company. 

    A 1:2 rights issue indicates that stockholders can buy 1 share for every 2 shares held. It's a right (not obligation) and can be traded. The share price also gets reduced in the same proportion. The number of issued shares also increases.

    \item \textbf{Share Buyback} is offered by a company to buy back its shares from the existing stockholders either because share price is too low or because it has surplus capital that it cannot put to good use that it plans to return to the shareholders. 

    It decreases the number of issued shares and increases EPS. One more reason can be because company is sure of its future growth. Consider buying if fundamentals are strong, revenue has not been decreased and will have enough liquidity thereafter to issue.

    \item \textbf{Delisting} of a company occurs either by a voluntary decision like below-par performance of the securities or a merger/acquisition of the listed company with another one or by the SEBI for violation of regulations. 

    It will require shareholder approval and buys them on a price bid by the shareholders, but it must be also acceptable to the company.

    \item \textbf{Priority of Payments Upon Company Sale:} When a company is sold, the proceeds from the sale are distributed in a specific order to fulfill the company's financial obligations. The typical priority of payments is as follows:
    \begin{enumerate}
      \item \textbf{Interest to Lenders:} Payments are first made to lenders, including those holding debt, bonds, and debentures.
      \item \textbf{Taxes to Government:} Next, tax liabilities owed to the government are settled.
      \item \textbf{Dividends to Preferred Shareholders:} If there are any preferred shareholders, they are paid their fixed dividends after taxes are cleared.
      \item \textbf{Reserves and Surplus to Common Equity Shareholders:} Finally, any remaining surplus is distributed to common equity shareholders, often referred to as retained earnings or reserves and surplus.
    \end{enumerate}

\end{itemize}

\end{document}